\documentclass[12pt,a4paper]{article}

% ============================================================================
% PACKAGES DE BASE
% ============================================================================
\usepackage{fontspec}
\usepackage[french]{babel}
\usepackage[left=2.5cm,right=2.5cm,top=2.5cm,bottom=2.5cm]{geometry}
\usepackage{setspace}
\usepackage{graphicx}
\usepackage{csquotes}
\usepackage{enumitem}
\usepackage{booktabs}
\usepackage{longtable}
\usepackage{multirow}
\usepackage{array}
\usepackage{xcolor}
\usepackage{adjustbox}
\usepackage{float}
\usepackage{hyperref}
\usepackage{fancyvrb}

% ============================================================================
% CONFIGURATION TYPOGRAPHIE
% ============================================================================
\IfFontExistsTF{Times New Roman}{%
  \setmainfont{Times New Roman}%
}{%
  \setmainfont{TeX Gyre Termes}%
}

\onehalfspacing
\graphicspath{{images/}}

% Pandoc helpers
\providecommand{\tightlist}{%
  \setlength{\itemsep}{0pt}\setlength{\parskip}{0pt}}
\providecommand{\pandocbounded}[1]{%
  \adjustbox{max width=0.94\linewidth,max totalheight=0.60\textheight,center}{#1}}
\setlength{\textfloatsep}{10pt plus 2pt minus 2pt}
\setlength{\intextsep}{10pt plus 2pt minus 2pt}
\newcommand{\AnnexeInput}[1]{%
  \VerbatimInput[fontsize=\small,frame=single]{#1}}

% ============================================================================
% HYPERLINKS ET MÉTADONNÉES PDF
% ============================================================================
\hypersetup{
  hidelinks,
  unicode=true,
  pdfauthor={Abdelatif DJEDDOU},
  pdftitle={Orchestration stigmergique de systèmes multi-agents LLM : principes de conception d'un framework adaptatif et application à la transformation de code},
  pdfsubject={Revue de littérature - Design Science Research - EMLV},
  pdfkeywords={stigmergie, multi-agents, LLM, Design Science Research, coordination organisationnelle, migration de code, gouvernance}
}

% ============================================================================
% MÉTADONNÉES DOCUMENT
% ============================================================================
\title{Orchestration stigmergique de systèmes multi-agents LLM : principes de conception d'un framework adaptatif et application à la transformation de code}
\author{Abdelatif DJEDDOU}
\date{6 février 2026}

% ============================================================================
% DÉBUT DU DOCUMENT
% ============================================================================
\begin{document}

\input{title_page_declaration.tex}
\tableofcontents
\newpage

\section{Introduction}\label{introduction}

En 2025, l\textquotesingle intelligence artificielle générative a
franchi un seuil décisif dans le domaine du génie logiciel : les grands
modèles de langage (LLM) ne se contentent plus de compléter du code
ligne par ligne, ils orchestrent désormais des workflows complets de
développement, de migration et de maintenance logicielle. Chez Google,
39 migrations de code distinctes totalisant 93 574 éditions
individuelles ont été réalisées avec l\textquotesingle assistance de LLM
en seulement douze mois, avec 74,45 \% des changements appliqués sans
modification humaine (Ziftci et al., 2025). Amazon Q Developer
transforme automatiquement des applications Java 8 vers Java 17 avec 79
\% du code appliqué directement (Amazon, 2024). IBM Watsonx
s\textquotesingle attaque à la modernisation de systèmes COBOL, estimant
que 775 milliards de lignes de code legacy restent en production à
l\textquotesingle échelle mondiale (IBM, 2023).

Ces avancées industrielles convergent vers un constat : la
transformation de code à grande échelle ne relève plus de la correction
unitaire de bugs, mais de la coordination de multiples agents logiciels
opérant simultanément sur des bases de code massives, interdépendantes
et critiques pour l\textquotesingle activité des entreprises.
Les organisations font face à des choix architecturaux complexes :
certaines cherchent à casser des monolithes pour migrer vers des
microservices, d\textquotesingle autres modernisent progressivement
des composants tout en maintenant le système principal en production,
d\textquotesingle autres encore doivent arbitrer entre refonte complète
et migration incrémentale. Cette diversité de stratégies exige des
cadres d\textquotesingle orchestration flexibles, capables de
s\textquotesingle adapter à des objectifs variés plutôt que de
prescrire un workflow unique.

Or, l\textquotesingle architecture dominante pour orchestrer ces agents
repose sur des modèles hiérarchiques où un superviseur central planifie,
distribue et contrôle les tâches de chaque agent subordonné, à
l\textquotesingle image des frameworks MetaGPT (Hong et al., 2024),
CrewAI ou AutoGen. Cette approche centralisée reproduit, dans le monde
logiciel, les structures organisationnelles classiques du
command-and-control.

Pourtant, la littérature récente révèle des limitations structurelles
préoccupantes de ces architectures hiérarchiques. Le framework MAST
identifie 14 modes d\textquotesingle échec distincts des systèmes
multi-agents LLM, avec des taux d\textquotesingle échec allant de 41 \%
à 86,7 \% sur sept frameworks état de l\textquotesingle art (Cemri et
al., 2025). L\textquotesingle approche Agentless démontre
qu\textquotesingle un pipeline séquentiel simple atteint 32 \% de
résolution sur SWE-bench Lite pour seulement 0,70 USD par problème,
surpassant tous les systèmes agents open source (Xia et al., 2024). Gao
et al. (2025) observent que les gains de performance des systèmes
multi-agents chutent de 10-16 \% à seulement 0,8-3 \% lorsque des
modèles frontières sont utilisés, questionnant la valeur ajoutée réelle
de la complexité agentique.

Face à ces constats, une alternative théorique existe : la stigmergie,
mécanisme de coordination indirecte où les traces laissées par une
action dans un médium stimulent les actions subséquentes (Heylighen,
2016a). Conceptualisée par Grassé (1959) pour expliquer la construction
collective des termitières sans superviseur central, la stigmergie a été
formalisée par Bonabeau, Dorigo et Theraulaz (1999) comme fondement de
l\textquotesingle intelligence en essaim, puis étendue par Heylighen
(2016a, 2016b) aux systèmes numériques contemporains tels que Wikipédia,
le développement open source et les systèmes de contrôle de version.
Ricci et al. (2007) ont proposé le cadre de la \emph{stigmergie
cognitive}, démontrant que des agents rationnels dotés de capacités
avancées peuvent se coordonner via des artefacts environnementaux
manipulables et instrumentables.

L\textquotesingle application de la stigmergie à
l\textquotesingle orchestration d\textquotesingle agents LLM pour la
transformation de code présente un intérêt théorique et pratique
considérable. D\textquotesingle un point de vue théorique, elle propose
de remplacer la communication explicite inter-agents, source de 36,94 \%
des échecs identifiés par MAST (Cemri et al., 2025), par une
coordination médiée par l\textquotesingle environnement.
D\textquotesingle un point de vue pratique, les systèmes de contrôle de
version comme Git fonctionnent déjà comme des médiums stigmergiques où
les commits stimulent les revues, tests et intégrations subséquentes
(Bolici et al., 2016). D\textquotesingle un point de vue managérial, la
traçabilité intrinsèque des artefacts partagés répond aux exigences
croissantes de gouvernance, d\textquotesingle auditabilité et de
conformité réglementaire, notamment celles de l\textquotesingle Article
14 de l\textquotesingle EU AI Act imposant une supervision humaine
effective des systèmes d\textquotesingle IA à haut risque (Fink, 2025).

Toutefois, la littérature académique présente un gap significatif : si
la stigmergie est bien théorisée en biologie et en systèmes distribués
classiques, et si les systèmes multi-agents LLM font
l\textquotesingle objet d\textquotesingle une recherche active, aucun
cadre intégrateur ne propose de combiner coordination stigmergique,
agents LLM et gouvernance organisationnelle dans un framework
généraliste applicable à des workflows complexes, dont la transformation
de code à grande échelle constitue un cas d\textquotesingle étude
privilégié. Les travaux existants traitent ces dimensions de
manière cloisonnée : la littérature sur la stigmergie ignore les LLM, la
littérature sur les agents LLM privilégie les architectures
hiérarchiques, et la littérature en management des systèmes
d\textquotesingle information aborde la gouvernance de
l\textquotesingle IA sans considérer les mécanismes de coordination
émergente. Plus fondamentalement, la théorie de la coordination
(Malone \& Crowston, 1994) identifie la stigmergie comme un mécanisme
organisationnel à part entière, au même titre que les hiérarchies et les
marchés, mais cette perspective n\textquotesingle a jamais été
opérationnalisée pour l\textquotesingle orchestration d\textquotesingle agents IA.

Ce constat motive la problématique suivante :

\begin{quote}
Comment concevoir un framework d\textquotesingle orchestration
multi-agents à coordination stigmergique, généraliste et gouvernable,
capable de surpasser les architectures hiérarchiques sur des tâches
complexes de coordination sous contraintes ?
La transformation de code constitue le domaine
d\textquotesingle application principal pour la validation empirique de
ces principes de conception.
\end{quote}

Cette problématique se décline en cinq objectifs de conception,
conformément à la méthodologie Design Science Research adoptée
(Hevner et al., 2004 ; Peffers et al., 2007) :

\begin{itemize}
\tightlist
\item
  OC1, Architecture généraliste : Concevoir un framework
  d\textquotesingle orchestration stigmergique (phéromones numériques,
  règles locales, environnement partagé) applicable à des tâches
  variées de coordination multi-agents, indépendamment du domaine.
\item
  OC2, Coordination émergente : Démontrer que les agents développent
  une spécialisation émergente et une allocation dynamique des tâches
  sans assignation explicite de rôles, conformément aux principes
  des systèmes adaptatifs complexes.
\item
  OC3, Supériorité sur la planification sous contraintes : Surpasser
  les frameworks centralisés sur le benchmark TravelPlanner (Xie et al.,
  2024) en score final et en satisfaction des contraintes dures.
\item
  OC4, Spécialisation à la transformation de code : Adapter le
  framework aux tâches de migration et de modernisation logicielle, et
  démontrer sa compétitivité sur les benchmarks PolyMigration (Amazon,
  2025) et SWE-bench (Jimenez et al., 2024) par rapport aux baselines
  existantes.
\item
  OC5, Gouvernance et auditabilité : Intégrer des mécanismes de
  traçabilité et de conformité réglementaire directement dans
  l\textquotesingle environnement stigmergique, et en valider la
  pertinence par évaluation d\textquotesingle experts.
\end{itemize}

La revue suit une progression continue : d\textquotesingle abord les
fondements de la stigmergie et leur ancrage dans les théories
managériales (coordination organisationnelle, capacités dynamiques,
routines organisationnelles), puis les mécanismes de coordination
computationnelle, les limites des architectures multi-agents, les usages
concrets en migration de code, et enfin les enjeux de gouvernance, de
processus et d\textquotesingle évaluation. L\textquotesingle objectif est
de construire un cadre cohérent, à la fois techniquement crédible et
ancré dans les sciences de gestion, conformément aux exigences du
Design Science Research en systèmes d\textquotesingle information
(Hevner et al., 2004).

\section{Fondements théoriques de la
stigmergie}\label{fondements-thuxe9oriques-de-la-stigmergie}

La stigmergie constitue le socle conceptuel de cette recherche. Pour
poser ce socle, on part de son origine en biologie, puis on suit son
passage vers les systèmes numériques contemporains.

\subsection{L\textquotesingle émergence du concept : Grassé et le
paradoxe de la
coordination}\label{luxe9mergence-du-concept--grassuxe9-et-le-paradoxe-de-la-coordination}

Le terme « stigmergie » fut introduit par le biologiste français
Pierre-Paul Grassé dans son étude pionnière des comportements de
construction chez les termites \emph{Bellicositermes natalensis} et
\emph{Cubitermes} sp. (Grassé, 1959). Grassé observa un paradoxe
fondamental : les termites semblaient travailler indépendamment les uns
des autres, sans coordination apparente, tout en produisant des
structures architecturales collectives d\textquotesingle une remarquable
complexité et cohérence. Ce « paradoxe de coordination » remettait en
question les deux théories dominantes de l\textquotesingle époque.
L\textquotesingle approche organiciste, inspirée de Spencer et Wheeler,
postulait l\textquotesingle existence de propriétés émergentes au niveau
de la colonie considérée comme un « superorganisme ».
L\textquotesingle approche individualiste, défendue par Rabaud,
soutenait que toute apparence de coopération n\textquotesingle était que
fortuite (Theraulaz \& Bonabeau, 1999).

Grassé résolut ce paradoxe en identifiant un troisième mécanisme : la
coordination indirecte médiée par l\textquotesingle environnement. Il
forgea le néologisme « stigmergie » à partir des racines grecques
\emph{stigma} (marque) et \emph{ergon} (travail), définissant ainsi la
stimulation des ouvriers par les performances qu\textquotesingle ils ont
eux-mêmes accomplies (Grassé, 1959). Le mécanisme stigmergique
fonctionne comme une boucle de rétroaction : une configuration
architecturale A déclenche une réponse comportementale R chez un
termite, qui modifie l\textquotesingle environnement pour créer une
nouvelle configuration A\textquotesingle, laquelle déclenche à son tour
une nouvelle réponse R\textquotesingle. Comme
l\textquotesingle expliquent Bonabeau et al. (1999, p. 18), chaque
travailleur crée de nouveaux stimuli en réponse aux configurations
existantes, ces stimuli agissant sur le même termite ou sur
n\textquotesingle importe quel autre membre de la colonie. Les
interactions pertinentes sont donc indirectes, transitant par
l\textquotesingle environnement.

\begin{figure}[H]
\centering
\pandocbounded{\includegraphics[keepaspectratio,alt={Boucle de rétroaction stigmergique : une action modifie l\textquotesingle environnement, ce qui déclenche l\textquotesingle action suivante. Source : adaptation à partir de Grassé (1959) et Bonabeau et al. (1999).}]{images/stigmergic_feedback_loop.png}}
\caption{Boucle de rétroaction stigmergique : une action modifie
l\textquotesingle environnement, ce qui déclenche
l\textquotesingle action suivante. Source : adaptation à partir de
Grassé (1959) et Bonabeau et al. (1999).}
\end{figure}

Cette découverte établit plusieurs principes fondamentaux pour la
compréhension des systèmes auto-organisés. Premièrement, la
décentralisation : aucun agent individuel ne possède une vision globale
ni n\textquotesingle exerce de contrôle centralisé. Deuxièmement, la
construction incrémentale : les agents exploitent ce que
d\textquotesingle autres ont construit pour apporter leur propre
contribution. Troisièmement, la flexibilité adaptative : lorsque
l\textquotesingle environnement change en raison d\textquotesingle une
perturbation externe, les agents répondent de manière appropriée sans
reprogrammation spécifique (Bonabeau et al., 1999). Quatrièmement, la
réduction de la communication : la stigmergie minimise le besoin de
communication directe entre agents, l\textquotesingle environnement
servant de mémoire partagée distribuée (Marsh \& Onof, 2008).

\begin{figure}[H]
\centering
\pandocbounded{\includegraphics[keepaspectratio,alt={Principes fondamentaux des systèmes auto-organisés : décentralisation, contribution incrémentale, adaptation et mémoire environnementale partagée. Source : Bonabeau et al. (1999) ; Marsh et Onof (2008).}]{images/fundamental_principles_self_organizing_systems.png}}
\caption{Principes fondamentaux des systèmes auto-organisés :
décentralisation, contribution incrémentale, adaptation et mémoire
environnementale partagée. Source : Bonabeau et al. (1999) ; Marsh et
Onof (2008).}
\end{figure}

\subsection{Formalisation algorithmique : l\textquotesingle intelligence
en essaim}\label{formalisation-algorithmique--lintelligence-en-essaim}

L\textquotesingle ouvrage fondateur \emph{Swarm Intelligence: From
Natural to Artificial Systems} (Bonabeau et al., 1999), publié dans la
série du Santa Fe Institute, établit l\textquotesingle intelligence en
essaim comme discipline scientifique légitime. Les auteurs y démontrent
systématiquement comment les principes d\textquotesingle autonomie, de
décentralisation et d\textquotesingle émergence peuvent remplacer le
contrôle centralisé traditionnel dans la résolution de problèmes
complexes. Leur contribution majeure réside dans
l\textquotesingle articulation du rôle de la modélisation comme
interface entre la compréhension de la nature et la conception de
systèmes artificiels (Bonabeau et al., 1999, p. 18). Cette approche
méthodologique permet d\textquotesingle explorer les modèles
biologiquement motivés au-delà de leurs contraintes naturelles initiales
, par exemple, les algorithmes d\textquotesingle optimisation par
colonie de fourmis (ACO) utilisent des taux de décroissance de
phéromones biologiquement implausibles mais computationnellement
efficaces.

Parallèlement, Theraulaz et Bonabeau (1999) clarifient dans leur
synthèse historique parue dans \emph{Artificial Life} la
conceptualisation de la stigmergie en distinguant deux mécanismes
fondamentaux. La stigmergie quantitative repose sur
l\textquotesingle intensité de la trace qui affecte la probabilité
d\textquotesingle une action : plus une piste de phéromone est
concentrée, plus la probabilité qu\textquotesingle une fourmi la suive
est élevée (Theraulaz \& Bonabeau, 1999, p. 104-105). La stigmergie
qualitative implique des stimuli qualitativement différents déclenchant
des réponses comportementales qualitativement différentes, comme dans la
construction de nids de guêpes où la configuration spatiale détermine le
type de réponse architecturale (Theraulaz \& Bonabeau, 1999, p. 105).
Cette typologie fournit un cadre conceptuel pour guider le choix entre
approches de coordination dans la conception de systèmes multi-agents.

\subsection{Extension aux systèmes numériques : la stigmergie comme
mécanisme
universel}\label{extension-aux-systuxe8mes-numuxe9riques--la-stigmergie-comme-muxe9canisme-universel}

Francis Heylighen entreprend, dans sa synthèse en deux parties publiée
dans \emph{Cognitive Systems Research}, la première formalisation
véritablement générale de la stigmergie (Heylighen, 2016a, 2016b). Il
propose une définition unifiée : la stigmergie est le mécanisme de
coordination indirecte où la trace laissée par une action dans un médium
stimule les actions subséquentes (Heylighen, 2016a). Cette définition
capture l\textquotesingle essence bidirectionnelle du processus :
l\textquotesingle action (\emph{ergon}) laisse une marque
(\emph{stigma}), qui stimule de nouvelles actions, formant une boucle de
rétroaction.

Heylighen (2016b) enrichit la typologie de Theraulaz et Bonabeau en
distinguant la stigmergie sématectonique, communication via modification
physique de l\textquotesingle environnement où la stimulation provient
du travail lui-même, et la stigmergie basée sur marqueurs, communication
via mécanismes de signalisation qui ne contribuent pas directement à la
tâche mais rendent les traces d\textquotesingle activité plus
saillantes. En contexte numérique, cette distinction
s\textquotesingle illustre avec Wikipédia : les lecteurs sont stimulés à
améliorer les pages soit directement en remarquant les lacunes du texte
(stigmergie sématectonique), soit indirectement via les notes « à faire
» résumant les tâches restantes (stigmergie basée sur marqueurs)
(Heylighen, 2016b, p. 51).

L\textquotesingle identification par Heylighen (2016a) de la stigmergie
comme mécanisme sous-jacent de multiples phénomènes de coordination
numérique est particulièrement pertinente pour cette recherche.
Wikipédia, le développement open source, les systèmes de contrôle de
version et même le mécanisme de prix du marché reposent sur des
dynamiques stigmergiques où les actions passées, réifiées dans un
médium, stimulent les actions futures. Heylighen (2016a, p. 6) explique
pourquoi ce mécanisme est resté longtemps sous-étudié :
l\textquotesingle interaction stigmergique étant par définition
indirecte, notre esprit biaisé pour rechercher des causes directes tend
à négliger la possibilité qu\textquotesingle un agent puisse influencer
le comportement d\textquotesingle un autre uniquement via la route
indirecte d\textquotesingle une trace laissée dans un environnement
passif.

\subsection{Épistémologie stigmergique et cognition
distribuée}\label{uxe9pistuxe9mologie-stigmergique-et-cognition-distribuuxe9e}

Marsh et Onof (2008) opèrent un pont entre stigmergie et philosophie de
l\textquotesingle esprit dans \emph{Cognitive Systems Research}. Partant
d\textquotesingle une conception située de la cognition, ils proposent
que la stigmergie constitue le mécanisme de coordination fondamental
pour comprendre comment la connaissance dans les communautés complexes
est essentiellement stigmergique (Marsh \& Onof, 2008, p. 136). Leur
caractérisation formelle des systèmes stigmergiques comporte trois
ensembles de propriétés : un environnement global comprenant des
environnements locaux avec une dynamique interne propre et une
perceptibilité seulement partielle ; des agents multiples présentant
rationalité limitée, comportement auto-organisé, stochastique et
adaptatif ; et des propriétés émergentes résultant des interactions,
propriétés qui ne sont ni prédictibles ni réductibles à des constituants
plus simples (Marsh \& Onof, 2008, p. 138).

Cette formalisation est directement applicable aux systèmes multi-agents
LLM : les agents disposent d\textquotesingle une rationalité limitée par
leur fenêtre de contexte, ils se comportent de manière stochastique
(génération probabiliste), et leurs interactions via un dépôt de code
partagé produisent des résultats collectifs non réductibles aux
contributions individuelles.

\subsection{Stigmergie cognitive et paradigme Agents \&
Artifacts}\label{stigmergie-cognitive-et-paradigme-agents--artifacts}

Ricci, Omicini, Viroli, Gardelli et Oliva (2007) proposent une extension
conceptuelle majeure avec la stigmergie cognitive, applicable aux agents
rationnels dotés de capacités cognitives avancées, par opposition aux
agents réactifs simples des approches biologiques classiques. Leur
cadre, publié dans \emph{Environments for Multi-Agent Systems III}
(Springer LNCS 4389), établit un pont fondamental entre stigmergie
classique et architectures de systèmes multi-agents (MAS) modernes.

Le paradigme Agents \& Artifacts repose sur une distinction
architecturale clé : les agents (entités computationnelles autonomes et
proactives) interagissent indirectement via des artefacts (ressources
environnementales manipulables supportant les activités des agents). Les
artefacts ne sont pas de simples objets passifs mais des médiateurs
actifs de coordination, fournissant des interfaces
d\textquotesingle interaction bien définies et rendant la coordination
indirecte explicitement instrumentable (Ricci et al., 2007).
Contrairement aux phéromones biologiques dont la dynamique est fixée par
la nature, les artefacts numériques peuvent être conçus, configurés et
adaptés pour répondre aux besoins spécifiques de coordination. Leur état
interne peut être exposé, permettant aux agents rationnels
d\textquotesingle observer les traces d\textquotesingle actions passées
et d\textquotesingle ajuster leur comportement. La durée de vie des
traces stigmergiques peut être contrôlée programmatiquement, offrant une
alternative à l\textquotesingle évaporation naturelle des phéromones.

Cette contribution s\textquotesingle avère fondamentale pour
l\textquotesingle orchestration multi-agents moderne basée sur les LLM,
car elle fournit le cadre conceptuel permettant de traiter les artefacts
numériques, documents partagés, bases de connaissances, systèmes de
tickets, logs d\textquotesingle exécution, comme médiateurs de
coordination stigmergique. Les agents LLM, dotés de capacités de
raisonnement complexes, peuvent interpréter et manipuler ces artefacts
avec une sophistication bien supérieure aux agents biologiques ou
réactifs traditionnels. En outre, les artefacts conservent
l\textquotesingle historique des modifications, permettant
l\textquotesingle audit a posteriori des décisions collectives, les
artefacts sont inspectables par des opérateurs humains maintenant le
\emph{human-in-the-loop}, et la coordination est externalisée dans des
artefacts dédiés séparant la logique métier de la mécanique de
coordination (Ricci et al., 2007).

\subsection{Systèmes adaptatifs complexes : Holland et
Kauffman}\label{systuxe8mes-adaptatifs-complexes--holland-et-kauffman}

Les travaux fondateurs de Holland (1995) et Kauffman (1993) fournissent
le cadre théorique des systèmes adaptatifs complexes (CAS) permettant de
comprendre comment l\textquotesingle organisation complexe émerge
d\textquotesingle agents adaptatifs simples. Holland (1995) introduit le
concept d\textquotesingle ordre caché pour décrire comment des patterns
organisationnels complexes peuvent émerger de règles
d\textquotesingle interaction simples entre agents, démontrant que la
complexité structurelle n\textquotesingle est pas nécessairement le
résultat d\textquotesingle une conception centralisée explicite.
Kauffman (1993) développe le concept d\textquotesingle ordre gratuit et
la dynamique au bord du chaos, montrant que les systèmes biologiques et
organisationnels opèrent dans une zone critique entre ordre rigide et
désordre chaotique où l\textquotesingle adaptabilité est maximale.

Ces fondations théoriques sont essentielles car elles établissent que
l\textquotesingle auto-organisation est une propriété fondamentale des
systèmes composés d\textquotesingle agents adaptatifs interagissant
localement, que la complexité peut émerger de la simplicité, et que les
systèmes opérant au bord du chaos peuvent s\textquotesingle adapter
dynamiquement aux perturbations sans reconfiguration centralisée. Ces
principes fournissent le fondement théorique pour concevoir des systèmes
d\textquotesingle orchestration multi-agents capables
d\textquotesingle auto-organisation, d\textquotesingle adaptation et de
coordination décentralisée.

Ces bases théoriques permettent maintenant de passer à la question
pratique : comment transformer cette logique de coordination indirecte
en mécanismes réellement utilisables dans des systèmes multi-agents
logiciels.

\subsection{Ancrage managérial : la stigmergie comme mécanisme de coordination organisationnelle}\label{ancrage-managuxe9rial--la-stigmergie-comme-muxe9canisme-de-coordination-organisationnelle}

Si les fondements biologiques et computationnels de la stigmergie sont
bien établis, son ancrage dans les théories du management et des
systèmes d\textquotesingle information reste insuffisamment exploité
dans la littérature sur les agents LLM. Trois cadres théoriques
permettent de repositionner la stigmergie comme mécanisme
organisationnel à part entière, dépassant ainsi la seule perspective
technique.

\subsubsection{Théorie de la coordination (Malone \& Crowston, 1994)}

Malone et Crowston (1994), dans un article fondateur publié dans
\emph{ACM Computing Surveys}, définissent la coordination comme la
gestion des dépendances entre activités. Leur taxonomie identifie
plusieurs mécanismes de coordination selon le type de dépendance :
les hiérarchies gèrent les dépendances par autorité centralisée, les
marchés par ajustement de prix, et la communication directe par
échanges de messages entre acteurs. La stigmergie constitue un
quatrième mécanisme, rarement identifié dans la littérature
managériale mais formellement cohérent avec le cadre de Malone et
Crowston : elle gère les dépendances par modification environnementale,
les agents se coordonnant en lisant et écrivant dans un artefact partagé
plutôt que par messages directs ou par soumission à une autorité
centrale.

Cette perspective est fondamentale car elle situe la stigmergie dans le
vocabulaire des sciences de gestion, non comme une curiosité biologique
transposée à l\textquotesingle informatique, mais comme un mécanisme de
coordination organisationnelle théorisé et comparable aux formes
classiques. Crowston lui-même a ultérieurement démontré empiriquement la
présence de coordination stigmergique dans les équipes de développement
open source distribuées (Bolici, Howison \& Crowston, 2016), validant la
transposition du concept au sein des organisations logicielles. Pour la
présente recherche, la théorie de la coordination fournit le cadre
analytique permettant d\textquotesingle évaluer quand et pourquoi la
coordination stigmergique surpasse les alternatives hiérarchiques ou
conversationnelles : elle réduit le nombre de liens de communication
de O($n^2$) à O($n$), élimine les goulots d\textquotesingle étranglement
centraux, et découple temporellement les activités interdépendantes.

\subsubsection{Capacités dynamiques et orchestration stratégique (Teece, 2007)}

Teece (2007), dans \emph{Strategic Management Journal}, conceptualise
les capacités dynamiques comme l\textquotesingle aptitude d\textquotesingle une
organisation à détecter les opportunités et menaces (\emph{sensing}),
à saisir les occasions (\emph{seizing}), et à reconfigurer ses actifs
pour maintenir un avantage compétitif (\emph{transforming}). Son concept
d\textquotesingle \emph{asset orchestration} désigne la capacité
managériale à assembler, structurer et recombiner les ressources
organisationnelles de manière dynamique.

Appliquée à la transformation de code, cette théorie requalifie
fondamentalement l\textquotesingle objet de recherche.
La modernisation logicielle continue n\textquotesingle est plus un
problème technique isolé mais une capacité dynamique que
l\textquotesingle organisation doit construire et maintenir : détecter
la dette technique et les opportunités de modernisation (\emph{sensing}),
mobiliser les ressources pour exécuter les transformations
(\emph{seizing}), et reconfigurer l\textquotesingle architecture
logicielle en réponse aux évolutions de l\textquotesingle environnement
technologique (\emph{transforming}). Le framework stigmergique proposé
dans ce mémoire opérationnalise cette capacité dynamique en fournissant
un mécanisme d\textquotesingle orchestration qui s\textquotesingle adapte
aux changements de contexte sans reconfiguration centralisée. Le
parallèle avec le concept d\textquotesingle \emph{asset orchestration}
de Teece est direct : les agents LLM orchestrés stigmergiquement
assemblent, transforment et reconfigurent les actifs logiciels de
l\textquotesingle organisation de manière distribuée et adaptative.

\subsubsection{Routines organisationnelles et évolution du code (Feldman \& Pentland, 2003)}

Feldman et Pentland (2003), dans un article influent publié dans
\emph{Administrative Science Quarterly}, reconceptualisent les routines
organisationnelles non comme des structures statiques mais comme des
sources de flexibilité et de changement. Ils distinguent
l\textquotesingle aspect ostensif (la structure abstraite de la routine
telle que comprise par les acteurs) et l\textquotesingle aspect
performatif (les actions concrètes réalisées par des personnes
spécifiques à des moments spécifiques). Les routines évoluent par
l\textquotesingle interaction continue entre ces deux aspects.

Cette théorie offre un éclairage puissant sur la transformation de code
vue comme processus organisationnel. Le code source incarne des routines
organisationnelles : chaque fonction, chaque module encode des
procédures, des règles métier et des pratiques qui structurent
l\textquotesingle activité de l\textquotesingle organisation. Transformer
du code, c\textquotesingle est donc transformer des routines
organisationnelles, avec toutes les implications de résistance, de
perte de savoir tacite et de recomposition des pratiques que cela
implique. La coordination stigmergique préserve l\textquotesingle aspect
ostensif des routines (les agents lisent et respectent la structure
existante du code comme contexte) tout en permettant l\textquotesingle évolution
de l\textquotesingle aspect performatif (les transformations concrètes
sont réalisées de manière incrémentale et traçable). Les traces
stigmergiques (commits, métadonnées, résultats de tests) constituent
la mémoire organisationnelle de l\textquotesingle évolution des
routines, permettant audit, apprentissage et réversibilité.

\subsubsection{Affordances technologiques (Strong et al., 2014)}

Strong et al. (2014) proposent dans \emph{MIS Quarterly} la théorie de
l\textquotesingle actualisation des affordances technologiques, qui
fournit un vocabulaire structuré pour analyser la valeur
organisationnelle des systèmes d\textquotesingle information. Une
affordance est une relation entre les propriétés d\textquotesingle un
objet technologique et les capacités d\textquotesingle un acteur
organisationnel qui permet une action potentielle. Appliqué au
framework stigmergique, ce cadre permet d\textquotesingle identifier
cinq affordances organisationnelles distinctes : la transformation
continue (les agents modifient le code de manière asynchrone et
incrémentale), l\textquotesingle auditabilité intrinsèque (les traces
environnementales créent automatiquement un historique complet),
la scalabilité adaptative (de nouveaux agents peuvent rejoindre la
colonie sans reconfiguration), la gouvernance distribuée (les
guardrails sont encodés dans l\textquotesingle environnement plutôt
que dans chaque agent), et la préservation des connaissances (les
artefacts partagés capitalisent le savoir collectif). Chacune de ces
affordances doit être actualisée par des pratiques organisationnelles
spécifiques, ce qui inscrit le déploiement du framework dans une
démarche de transformation organisationnelle et non de simple
implémentation technique.

Ces quatre cadres théoriques, théorie de la coordination, capacités
dynamiques, routines organisationnelles et affordances technologiques,
constituent l\textquotesingle ancrage managérial de ce mémoire. Ils
permettent de répondre à l\textquotesingle exigence d\textquotesingle un
mémoire en management des systèmes d\textquotesingle information : la
contribution ne réside pas uniquement dans l\textquotesingle artefact
technique, mais dans les principes de conception
généralisables qui émergent de
l\textquotesingle intersection entre théorie organisationnelle et
ingénierie des systèmes multi-agents.

\section{Coordination multi-agents et systèmes
distribués}\label{coordination-multi-agents-et-systuxe8mes-distribuuxe9s}

La transposition des principes stigmergiques vers des systèmes
informatiques distribués demande des mécanismes concrets. On clarifie
d\textquotesingle abord les notions d\textquotesingle agent et de
système multi-agent, puis on regarde comment la coordination passe par
un environnement partagé.

\subsection{Agents LLM et systèmes multi-agents : définitions
opératoires}\label{agents-llm-et-systuxe8mes-multi-agents--duxe9finitions-opuxe9ratoires}

Avant de détailler la coordination, il faut distinguer deux niveaux
souvent mélangés dans les publications récentes :
l\textquotesingle agent LLM et le système multi-agent (Cemri et al.,
2025). Un agent LLM est un logiciel qui reçoit des consignes, garde un
état (historique ou mémoire) et agit dans un environnement via des
outils. Un système single-agent concentre tout dans une seule boucle
d\textquotesingle action. Un système multi-agent répartit le travail
entre plusieurs agents spécialisés qui se coordonnent entre eux (Cemri
et al., 2025 ; Gao et al., 2025).

La littérature avance des bénéfices clairs pour les MAS : découpage des
tâches, parallélisation et spécialisation des rôles. Mais ce choix a
aussi des coûts : plus de latence, plus de tokens, plus de difficulté de
débogage et de gouvernance (Gao et al., 2025 ; Kapoor et al., 2024).

Un constat important pour cette recherche est que les MAS LLM dominants
s\textquotesingle appuient sur des patterns de coordination
conversationnels (échanges de messages) ou hiérarchiques
(superviseur-workers). Rodriguez (2026) formalise une alternative via
une coordination par champs de pression et décroissance temporelle, où
les agents opèrent localement sur un artefact commun. Ce déplacement
vers l\textquotesingle environnement comme support de coordination est
cohérent avec la perspective \emph{environment-first} de Weyns et al.
(2007), qui conceptualisent l\textquotesingle environnement comme une
abstraction de première classe dans les systèmes multi-agents.

\begin{figure}[H]
\centering
\pandocbounded{\includegraphics[keepaspectratio,alt={Agent LLM vs système multi-agent (MAS) : comparaison des niveaux d\textquotesingle organisation et des modes de coordination. Source : adaptation de Cemri et al. (2025) et Gao et al. (2025).}]{images/agentLLMvsMAS.png}}
\caption{Agent LLM vs système multi-agent (MAS) : comparaison des
niveaux d\textquotesingle organisation et des modes de coordination.
Source : adaptation de Cemri et al. (2025) et Gao et al. (2025).}
\end{figure}

\subsection{Phéromones numériques : de la biologie au
calcul}\label{phuxe9romones-numuxe9riques--de-la-biologie-au-calcul}

Parunak (1997) développe dans son article fondateur une traduction
systématique des principes biologiques en principes
d\textquotesingle ingénierie pour systèmes multi-agents, identifiant
trois principes fondamentaux : localité de perception et
d\textquotesingle action, coordination stigmergique, et feedback
positif. Il souligne que les agents dans les systèmes naturels sont
remarquablement simples, mais que leur interaction produit des
comportements collectifs sophistiqués (Parunak, 1997, p. 70). Cette
simplicité individuelle couplée à la complexité émergente constitue un
avantage architectural majeur : la robustesse du système dépend de la
redondance et de l\textquotesingle adaptabilité collective plutôt que de
la sophistication de composants individuels fragiles.

\subsection{Espaces de tuples et communication
générative}\label{espaces-de-tuples-et-communication-guxe9nuxe9rative}

Le modèle Linda, introduit par Gelernter (1985), révolutionne la
coordination dans les systèmes distribués en proposant la communication
générative via un espace de tuples, mémoire associative partagée où les
processus déposent et récupèrent des tuples sans connaître
l\textquotesingle identité des autres processus. Gelernter (1985, p. 81)
critique le modèle classique de message-passing qui crée un couplage
temporel et spatial rendant les programmes fragiles. Linda permet une
coordination découplée selon trois dimensions : temporelle (producteur
et consommateur n\textquotesingle ont pas besoin
d\textquotesingle exister simultanément), spatiale (les processus ne
connaissent pas l\textquotesingle identité les uns des autres), et de
synchronisation (la communication est asynchrone).

L\textquotesingle espace de tuples constitue une concrétisation
computationnelle du médium stigmergique théorisé par Heylighen (2016a).
Les tuples représentent les traces laissées par les actions des agents,
qui stimulent les actions subséquentes d\textquotesingle autres agents.
Gelernter (1985, p. 89) observe que les données acquièrent une existence
indépendante dans l\textquotesingle espace de tuples, incarnant le
principe stigmergique de médiation environnementale.

Carriero et Gelernter (1992) formalisent un principe architectural
fondamental : la séparation entre modèle de calcul et modèle de
coordination. La coordination doit être traitée comme une préoccupation
de première classe, distincte et orthogonale au calcul (Carriero \&
Gelernter, 1992, p. 98). Cette orthogonalité justifie théoriquement la
séparation entre services métier et mécanismes
d\textquotesingle orchestration, fondement des architectures de
microservices contemporaines.

Mamei et Zambonelli (2009) proposent l\textquotesingle extension TOTA
(\emph{Tuples On The Air}) introduisant des tuples distribués capables
de se propager de manière autonome à travers le réseau, de
s\textquotesingle auto-maintenir face aux changements topologiques, et
de se transformer selon leur contexte de propagation. Cette approche
génère des champs computationnels, gradients distribués
d\textquotesingle information guidant les agents, et
s\textquotesingle avère particulièrement pertinente pour les
architectures de microservices dynamiques (Mamei \& Zambonelli, 2009).

\subsection{Auto-organisation dans les systèmes
multi-agents}\label{auto-organisation-dans-les-systuxe8mes-multi-agents}

Serugendo, Gleizes et Karageorgos (2005) proposent une taxonomie
complète des mécanismes d\textquotesingle auto-organisation, définissant
l\textquotesingle auto-organisation comme un processus dynamique et
adaptatif où les systèmes acquièrent et maintiennent une structure sans
contrôle externe explicite (Serugendo et al., 2005, p. 166). Weyns,
Omicini et Odell (2007) complètent cette perspective en proposant
l\textquotesingle environnement comme abstraction de première classe
dans les systèmes multi-agents, argumentant que les environnements ne
sont pas de simples conteneurs passifs mais des entités actives
structurant la coordination. Ce positionnement résonne directement avec
le cadre stigmergique où l\textquotesingle environnement est le médium
de coordination plutôt qu\textquotesingle un simple support
d\textquotesingle exécution.

\subsection{Stigmergie empirique dans le développement open
source}\label{stigmergie-empirique-dans-le-duxe9veloppement-open-source}

Bolici, Howison et Crowston (2016) fournissent une validation empirique
de la coordination stigmergique dans le développement logiciel libre
(FLOSS). Leur étude, publiée dans \emph{Cognitive Systems Research},
démontre que les équipes FLOSS utilisent des mécanismes stigmergiques
via les systèmes de contrôle de version, les systèmes de suivi de bugs,
et les listes de diffusion, où les contributions de code servent de
stimuli déclenchant les activités subséquentes de revue, test et
intégration (Bolici et al., 2016). La traçabilité intrinsèque des
systèmes de contrôle de version crée automatiquement une trace complète
et auditable des modifications.

Cette observation est directement pertinente pour notre problématique :
Git, outil omniprésent du génie logiciel, fonctionne déjà comme un
médium stigmergique. Les commits représentent les traces, les branches
les chemins de phéromones, les pull requests les points de validation
collective, et les tests CI/CD les mécanismes de feedback
environnemental. L\textquotesingle orchestration multi-agents
stigmergique ne nécessite donc pas la création d\textquotesingle une
infrastructure nouvelle, mais l\textquotesingle instrumentation et
l\textquotesingle extension d\textquotesingle outils déjà adoptés
industriellement.

\subsection{Développements récents : bio-inspiration dans
l\textquotesingle orchestration agentique
(2023-2026)}\label{duxe9veloppements-ruxe9cents--bio-inspiration-dans-lorchestration-agentique-2023-2026}

Un développement notable du SOTA 2023-2026 est la réintroduction
explicite de mécanismes d\textquotesingle intelligence en essaim dans
l\textquotesingle orchestration agentique. Chari et al. (2025) proposent
ACO-ToT, combinant LLM et optimisation par colonie de fourmis pour
apprendre des chemins de raisonnement via renforcement de phéromones
dans un arbre de pensée. Zhang et al. (2025) introduisent SwarmAgentic,
un cadre visant à automatiser la conception de systèmes multi-agents via
une population de systèmes candidats évoluant par mises à jour guidées
par feedback, inspirées du Particle Swarm Optimization. Rodriguez (2026)
formalise une coordination émergente par champs de pression et
décroissance temporelle, où les agents opèrent localement sur un
artefact commun.

SwarmAgentic mérite une attention particulière car il démontre la
viabilité d\textquotesingle un framework généraliste inspiré de
l\textquotesingle intelligence en essaim. Évalué sur le benchmark
TravelPlanner (Xie et al., 2024), qui mesure la capacité des systèmes
agents à planifier des itinéraires sous contraintes multiples
(budget, durée, préférences), SwarmAgentic obtient avec GPT-4o un score
final de 32,2\%, surpassant nettement les baselines : Direct (2,2\%),
Chain-of-Thought (2,2\%), Self-Refine (1,1\%), EvoAgent (7,2\%) et ADAS
(8,9\%) (Zhang et al., 2025). Le score de satisfaction des contraintes
dures (\emph{Hard Constraint Macro}) atteint 52,8\% avec GPT-4o, contre
27,8\% pour ADAS, le second meilleur système. Ces résultats sont
significatifs pour deux raisons : ils montrent qu\textquotesingle une
coordination bio-inspirée peut surpasser des approches
conversationnelles classiques sur des tâches de planification complexe,
et ils établissent TravelPlanner comme benchmark de référence pour
l\textquotesingle évaluation de frameworks d\textquotesingle orchestration
généralistes, au-delà du seul génie logiciel.

Ces travaux convergent vers le positionnement de cette recherche : le
déplacement d\textquotesingle une partie de la coordination depuis des
protocoles conversationnels explicites vers un médium partagé et des
signaux d\textquotesingle évaluation, favorisant
l\textquotesingle adaptation et réduisant l\textquotesingle overhead
d\textquotesingle orchestration. Rahman et al. (2025) interrogent
toutefois la transposition du terme « swarm » aux frameworks LLM,
soulignant les différences fondamentales entre agents biologiques
simples et agents LLM complexes dotés de capacités de raisonnement
avancées. Cette tension justifie le recours au concept de
\emph{stigmergie cognitive} (Ricci et al., 2007) plutôt qu\textquotesingle à
l\textquotesingle intelligence en essaim classique : les agents LLM ne
sont pas des fourmis réactives mais des agents cognitifs capables
d\textquotesingle interpréter et de manipuler des artefacts
environnementaux complexes.

\subsection{Panorama des frameworks et
benchmarks}\label{panorama-des-frameworks-et-benchmarks}

Deux familles complémentaires de frameworks coexistent pour concevoir
des systèmes multi-agents : les plateformes MAS classiques pré-LLM
(JADE, Jason, SPADE) reposant sur le message-passing et des modèles
organisationnels explicites, et les frameworks LLM-native (LangGraph,
AutoGen, CrewAI, MetaGPT) orientés orchestration
d\textquotesingle agents LLM intégrant tool-use, gestion
d\textquotesingle état et topologies d\textquotesingle interaction (Geng
\& Chang, 2025 ; Cemri et al., 2025). MetaGPT formalise des
\emph{Standardized Operating Procedures} (SOPs) pour réduire les
incohérences des chaînes d\textquotesingle agents dans les tâches de
génie logiciel (Hong et al., 2024). OpenHands propose une plateforme
open source pour développer et évaluer des agents généralistes capables
d\textquotesingle interagir avec dépôts, shell et tests (Wang et al.,
2025).

L\textquotesingle essor des agents s\textquotesingle accompagne de
benchmarks dédiés : AgentBench vise l\textquotesingle évaluation «
LLM-as-Agent » dans plusieurs environnements (Liu et al., 2024),
SWE-bench structure l\textquotesingle évaluation sur des issues GitHub
réelles (Jimenez et al., 2024), et SWE-agent met en avant
l\textquotesingle importance des interfaces agent-ordinateur (Yang et
al., 2024). Ces benchmarks sont repris plus loin dans la section dédiée
à l\textquotesingle évaluation.

Ce cadrage permet d\textquotesingle aborder le point suivant de manière
plus lucide : même si les MAS promettent beaucoup, ils ont aussi des
limites concrètes qu\textquotesingle il faut intégrer dès la conception.

\section{Limites des approches
multi-agents}\label{limites-des-approches-multi-agents}

Pour garder une lecture équilibrée, il faut aussi regarder les limites
des architectures multi-agents et les cas où des approches plus simples
fonctionnent mieux.

\subsection{L\textquotesingle approche Agentless : la simplicité contre
la
complexité}\label{lapproche-agentless--la-simplicituxe9-contre-la-complexituxe9}

L\textquotesingle article Agentless (Xia et al., 2024) constitue une
contribution majeure remettant en question les présupposés de
l\textquotesingle ingénierie logicielle basée sur des agents autonomes.
Xia et al. (2024) proposent une approche simplifiée structurée en trois
phases séquentielles : localisation hiérarchique des emplacements
d\textquotesingle édition pertinents, génération de multiples patches
candidats en format diff, et validation via tests de reproduction et
régression. Contrairement aux outils basés sur agents, Agentless
contient des étapes bien définies sans laisser l\textquotesingle agent
LLM décider des actions futures ni utiliser d\textquotesingle outils
complexes (Xia et al., 2024, p. 4).

Les résultats sur SWE-bench Lite sont remarquables : 32 \% de précision
(96 corrections sur 300 problèmes) pour 0,70 USD de coût moyen par
problème, surpassant tous les systèmes agents open source (Xia et al.,
2024, p. 3). Cette performance est d\textquotesingle autant plus
significative qu\textquotesingle elle est obtenue sans mécanismes
complexes d\textquotesingle agents autonomes. Les auteurs espèrent
qu\textquotesingle Agentless aidera à réinitialiser le baseline et le
point de départ pour les agents logiciels autonomes (Xia et al., 2024,
p. 3).

Il convient de contextualiser cette critique : Agentless est focalisé
sur le software engineering avec SWE-bench comme benchmark spécifique.
La généralisation à des workflows de migration multi-fichiers à grande
échelle, où les dépendances inter-composants créent une complexité
structurelle supérieure, nécessite une validation empirique
supplémentaire.

\subsection{Framework MAST : comprendre où les MAS
échouent}\label{framework-mast--comprendre-ouxf9-les-mas-uxe9chouent}

Le framework MAST (Cemri et al., 2025) est utile parce
qu\textquotesingle il montre clairement les familles
d\textquotesingle échecs récurrents dans les systèmes multi-agents LLM.
En pratique, ces échecs reviennent souvent à trois problèmes : des
tâches mal cadrées, des agents qui se comprennent mal entre eux, et des
sorties insuffisamment vérifiées avant la fin du processus.

Même sans entrer dans les détails méthodologiques, le message est fort :
les taux d\textquotesingle échec observés restent élevés sur plusieurs
frameworks open source, ce qui indique un enjeu de conception, pas
seulement de puissance de modèle. Pour un lecteur orienté management SI,
cela renforce l\textquotesingle idée qu\textquotesingle un bon système
ne dépend pas uniquement du modèle, mais surtout de
l\textquotesingle architecture de coordination et des mécanismes de
contrôle.

\begin{figure}[H]
\centering
\pandocbounded{\includegraphics[keepaspectratio,alt={Taxonomie MAST des modes d\textquotesingle échec des systèmes multi-agents LLM. Source : Cemri et al. (2025).}]{images/taxomanie_MAST.png}}
\caption{Taxonomie MAST des modes d\textquotesingle échec des systèmes
multi-agents LLM. Source : Cemri et al. (2025).}
\end{figure}

\subsection{Analyse coût-efficacité et érosion de
l\textquotesingle avantage
multi-agents}\label{analyse-couxfbt-efficacituxe9-et-uxe9rosion-de-lavantage-multi-agents}

Kapoor et al. (2024), dans leur article « AI Agents That Matter »
(Princeton University), démontrent que l\textquotesingle optimisation
mono-objectif (précision seule) masque des inefficacités systémiques
majeures. Leur analyse révèle des variations de coût allant
jusqu\textquotesingle à 50x entre agents atteignant des niveaux de
précision similaires, et introduisent le concept de frontière de Pareto
coût-précision comme méthodologie d\textquotesingle évaluation
rigoureuse. La précision seule ne peut identifier le progrès car elle
peut être améliorée par des méthodes scientifiquement dénuées de sens
telles que le retry (Kapoor et al., 2024, p. 4).

Gao et al. (2025) approfondissent cette critique en comparant
systématiquement systèmes single-agent et multi-agents. Leurs résultats
montrent que les gains de performance MAS chutent de 10-16 \% avec des
modèles de capacité moyenne à seulement 0,8-3 \% avec des modèles
frontières, pour un surcoût en tokens de 4x à 220x (Gao et al., 2025).
Les auteurs proposent des paradigmes hybrides (routing, cascade)
améliorant simultanément précision (+1,1-12 \%) et efficacité (-88,1 \%
de coût), suggérant que la question n\textquotesingle est pas «
single-agent ou multi-agent ? » mais comment adapter dynamiquement le
degré de coordination à la complexité de la tâche.

\subsection{Positions industrielles et principe de
parcimonie}\label{positions-industrielles-et-principe-de-parcimonie}

Yan (2025), CTO de Cognition AI (créateurs de Devin), argumente dans un
billet de blog technique que la fragmentation de contexte constitue le
défaut architectural fondamental des systèmes multi-agents. Lorsque le
contexte est distribué entre multiples agents, aucun agent individuel ne
possède une vue suffisamment complète pour prendre des décisions
optimales.

Anthropic (2025) recommande officiellement de commencer avec
l\textquotesingle approche la plus simple qui fonctionne et
d\textquotesingle ajouter de la complexité seulement quand
l\textquotesingle évidence le supporte. Cette guidance méthodologique
s\textquotesingle inscrit dans la tradition du rasoir
d\textquotesingle Occam appliqué à l\textquotesingle ingénierie
logicielle.

Cette lecture critique ne disqualifie pas la stigmergie ; elle fixe
plutôt un niveau d\textquotesingle exigence. La suite doit donc montrer,
sur des cas de migration concrets, qu\textquotesingle une coordination
par environnement peut réellement améliorer le compromis
performance-coût-gouvernabilité.

\section{Migration de code et transformation avec les
LLM}\label{migration-de-code-et-transformation-avec-les-llm}

La migration et la transformation de code constituent le terrain
d\textquotesingle application principal de cette recherche. Ce terrain
permet de relier directement les choix d\textquotesingle orchestration à
des contraintes industrielles réelles : qualité, coût, délai et maîtrise
du risque.

\subsection{Migration à grande échelle : la validation industrielle de
Google}\label{migration-uxe0-grande-uxe9chelle--la-validation-industrielle-de-google}

L\textquotesingle article de Ziftci et al. (2025), présenté au FSE 2025
(conférence top-tier en génie logiciel), constitue la première
validation industrielle académiquement documentée de migration
LLM-assistée à très grande échelle. Sur 39 migrations distinctes
totalisant 595 changements de code et 93 574 éditions individuelles,
74,45 \% des changements générés ont été appliqués sans modification
humaine, réduisant de 50 \% le temps développeur total requis (Ziftci et
al., 2025, p. 6). L\textquotesingle architecture repose sur une
combinaison d\textquotesingle analyse statique pour
l\textquotesingle identification des emplacements, de génération LLM
pour les patches, de validation automatisée via tests existants, et de
review humaine sélective pour les cas complexes.

Malgré ces résultats, Ziftci et al. (2025) identifient des limitations
significatives. La fenêtre de contexte limitée des LLM contraint
l\textquotesingle efficacité pour les fichiers volumineux ou les
migrations nécessitant la compréhension de multiples fichiers
interdépendants (p. 11). Les tests pré-existants défaillants
(\emph{flaky tests}) empêchent la validation automatique (p. 12).
Environ 25 \% des changements nécessitent une intervention humaine pour
des cas non anticipés.

Ces principes sont directement transposables à
l\textquotesingle orchestration multi-agents stigmergique : les tests
automatisés servent de feedback stigmergique validant la qualité des
modifications, les changements peuvent être générés, validés et intégrés
de manière asynchrone, et chaque changement est tracé dans le système de
contrôle de version fournissant l\textquotesingle auditabilité requise
pour la gouvernance.

\subsection{Coordination multi-agents pour le codage autonome : les
leçons de
Cursor}\label{coordination-multi-agents-pour-le-codage-autonome--les-leuxe7ons-de-cursor}

L\textquotesingle article de Cursor (2025) documente les leçons tirées
d\textquotesingle expérimentations avec des agents de codage autonomes
fonctionnant pendant des semaines consécutives.
L\textquotesingle approche initiale reposait sur une coordination
décentralisée égalitaire où chaque agent vérifiait ce que les autres
faisaient via un fichier partagé avec mécanismes de verrouillage. Cette
approche a échoué de manière instructive : goulots
d\textquotesingle étranglement des verrous (vingt agents ralentissant au
débit effectif de deux ou trois), fragilité systémique (défaillances en
cascade), et aversion au risque collective (sans hiérarchie, les agents
évitaient les tâches difficiles) (Cursor, 2025).

Le passage à un contrôle de concurrence optimiste puis à une
architecture planificateurs-travailleurs illustre une trajectoire
pertinente pour la conception d\textquotesingle orchestrations
stigmergiques. L\textquotesingle échec de la coordination décentralisée
naïve confirme que la stigmergie ne peut se réduire à un simple fichier
partagé : elle nécessite des mécanismes de sélection, de renforcement et
d\textquotesingle évaporation inspirés des systèmes biologiques pour
produire une coordination effective.

\subsection{Traduction cross-language et modernisation
legacy}\label{traduction-cross-language-et-modernisation-legacy}

Rozière et al. (2020), dans leur article présenté à NeurIPS 2020,
démontrent la possibilité de traduction non supervisée entre langages de
programmation en exploitant les mécanismes de \emph{back-translation} et
de \emph{denoising auto-encoding}, atteignant des scores BLEU de 68,7
pour la traduction C++ vers Java et 61,9 pour Java vers C++ sans données
d\textquotesingle entraînement parallèles (le score BLEU mesure la
proximité entre une traduction générée et une traduction de référence :
plus il est élevé, plus les deux sont proches). CodeRosetta (2024) étend
cette approche à la programmation parallèle (NeurIPS 2024). Ces travaux
établissent que la traduction cross-language par LLM est techniquement
faisable, ouvrant la voie à son orchestration dans des workflows de
migration multi-fichiers.

La modernisation de systèmes COBOL vers Java représente un cas
industriel critique. Sneed (2010) documente les défis fondamentaux de
cette migration : différences paradigmatiques entre COBOL procédural et
Java orienté-objet, complexité des systèmes legacy (775 milliards de
lignes en production selon IBM, 2023), et nécessité de préserver la
sémantique métier exacte. Cet enjeu est central pour les secteurs
bancaire, assurance et paiements, où une partie du cœur opérationnel
repose encore sur du COBOL. IBM Watsonx combine analyse statique et
génération LLM pour automatiser cette transformation (IBM, 2023). Diggs
et al. (2025) évaluent l\textquotesingle utilisation de LLM pour générer
la documentation de code legacy comme prérequis à la modernisation,
publié dans le cadre du workshop LLM4Code@ICSE 2025.

\subsection{Évolution et refactoring
d\textquotesingle API}\label{uxe9volution-et-refactoring-dapi}

Lamothe, Guéhéneuc et Shang (2021) produisent dans \emph{ACM Computing
Surveys} la revue systématique la plus complète sur
l\textquotesingle évolution des API, analysant 110 études primaires.
Leur synthèse identifie les ruptures d\textquotesingle API comme un
problème central du génie logiciel et cartographie les solutions
existantes pour la migration de clients lors
d\textquotesingle évolutions d\textquotesingle API. Dig et Johnson
(2005) classifient les refactorings de manière à orienter les outils de
migration automatisés. Ces travaux fondent les exigences fonctionnelles
pour un système multi-agents de migration : identification automatique
des breaking changes, génération de patches de mise à jour, et
validation de la préservation sémantique.

\subsection{Transformation du développeur par
l\textquotesingle IA}\label{transformation-du-duxe9veloppeur-par-lia}

Peng et al. (2023), dans une étude contrôlée randomisée, documentent une
amélioration de 55,8 \% de la productivité des développeurs utilisant
GitHub Copilot. Barke et al. (2023) montrent de leur côté que le travail
alterne entre un mode « accélération » (l\textquotesingle IA aide à
produire plus vite) et un mode « exploration » (l\textquotesingle IA
aide à clarifier un problème encore flou).

Ce déplacement du rôle apparaît encore plus clairement dans des retours
récents de terrain. Le billet OpenAI du 5 février 2026 sur GPT‑5.3‑Codex
décrit des équipes qui utilisent Codex pour accélérer leur propre cycle
de recherche et de déploiement, avec des chercheurs et ingénieurs
expliquant que leur travail est déjà « fondamentalement différent » de
celui d\textquotesingle il y a deux mois (OpenAI, 2026). Côté Anthropic,
une étude interne publiée en décembre 2025 montre aussi un glissement
vers des tâches de pilotage, supervision et arbitrage, plutôt
qu\textquotesingle une production ligne par ligne en continu (Anthropic,
2025b).

Dans ce contexte, la \emph{pull request} (PR, proposition de changement
de code soumise à revue) devient un objet clé : le développeur
n\textquotesingle est plus seulement auteur direct du code, il devient
aussi superviseur de sorties générées, arbitre de qualité et garant de
cohérence.

Ces retours terrain préparent logiquement la suite : si les agents
peuvent produire beaucoup plus vite, la question centrale devient la
gouvernance de cette production à l\textquotesingle échelle.

\section{Gouvernance, auditabilité et
conformité}\label{gouvernance-auditabilituxe9-et-conformituxe9}

La gouvernance devient ici le point d\textquotesingle équilibre entre
autonomie des agents et contrôle organisationnel.

\subsection{Guardrails pour écologies
humain-IA}\label{guardrails-pour-uxe9cologies-humain-ia}

Grisold, Berente et Seidel (2025) publient dans \emph{MIS Quarterly} une
contribution théorique majeure introduisant le concept
d\textquotesingle écologies humain-IA, définies comme des environnements
où agents humains et IA se coordonnent via des normes sociales plutôt
que par programmation explicite (Grisold et al., 2025, p. 1241). Leur
concept central de \emph{guardrails} est défini comme des zones de
comportement désirable qui contraignent les opérations système sans
éliminer l\textquotesingle autonomie adaptative (Grisold et al., 2025,
p. 1243).

Les auteurs distinguent les normes profondes (\emph{deep norms}),
principes fondamentaux stables reflétant les valeurs organisationnelles
de long terme, et les normes de surface (\emph{surface norms}),
guardrails opérationnels spécifiques pouvant évoluer selon les
contextes. Leur Design Theory for Predictability
s\textquotesingle articule autour de trois principes :
l\textquotesingle explicitation sélective des normes, la transparence
des frontières comportementales, et l\textquotesingle évolutivité des
guardrails sans reconfiguration système complète (Grisold et al., 2025,
p. 1250-1253).

Cette approche possède une pertinence directe pour
l\textquotesingle orchestration multi-agents stigmergique. Dans un
système stigmergique, les guardrails peuvent être implémentés comme
contraintes environnementales plutôt que comme règles programmatiques
dans chaque agent. L\textquotesingle environnement (dépôt Git) peut
n\textquotesingle accepter que les commits signés et liés à une issue
tracker (guardrail structurel), déclencher automatiquement un workflow
de revue humaine au-delà d\textquotesingle un seuil de complexité
(guardrail de validation), et révoquer automatiquement toute
modification échouant aux tests CI/CD (guardrail de rollback). Cette
approche centralise la gouvernance dans l\textquotesingle environnement
plutôt que de la distribuer dans chaque agent, facilitant
l\textquotesingle auditabilité et l\textquotesingle évolution des
politiques.

\subsection{Perspective principal-agent et
orchestration}\label{perspective-principal-agent-et-orchestration}

Jarrahi et Ritala (2025), dans \emph{California Management Review},
appliquent la théorie économique principal-agent aux agents IA,
démontrant que les agents IA travaillent au nom des utilisateurs dans
des contraintes définies mais disposent d\textquotesingle une capacité
de traitement et d\textquotesingle accès à l\textquotesingle information
supérieurs (Jarrahi \& Ritala, 2025, p. 2). Leur framework à cinq
dimensions pour la gouvernance, interactivité et adaptabilité,
autonomie guidée, intelligence collective, sécurité, responsabilité et
interopérabilité via orchestration, individualisation et alignement,
fournit un cadre structurant pour la conception de systèmes multi-agents
gouvernables (Jarrahi \& Ritala, 2025).

\begin{figure}[H]
\centering
\pandocbounded{\includegraphics[keepaspectratio,alt={Cadre de gouvernance pour systèmes multi-agents : articulation des dimensions de contrôle, responsabilité et alignement. Source : adaptation de Jarrahi et Ritala (2025).}]{images/cadre_de_gouvernance.png}}
\caption{Cadre de gouvernance pour systèmes multi-agents : articulation
des dimensions de contrôle, responsabilité et alignement. Source :
adaptation de Jarrahi et Ritala (2025).}
\end{figure}

L\textquotesingle emphase sur l\textquotesingle auditabilité et la
traçabilité des décisions agents comme prérequis pour
l\textquotesingle accountability organisationnelle (Jarrahi \& Ritala,
2025, p. 9) converge avec l\textquotesingle approche stigmergique via
systèmes de contrôle de version, où l\textquotesingle historique des
modifications reste inspectable.

\subsection{Conformité réglementaire : EU AI Act Article
14}\label{conformituxe9-ruxe9glementaire--eu-ai-act-article-14}

L\textquotesingle Article 14 de l\textquotesingle EU AI Act impose que
les systèmes d\textquotesingle IA à haut risque soient conçus et
développés de manière à permettre une supervision humaine effective
durant leur utilisation (cité dans Fink, 2025, p. 3). Cette formulation
établit deux principes : le \emph{design for oversight} (la supervision
doit être intégrée dès la conception) et
l\textquotesingle{}\emph{effectiveness} (la supervision doit être
significative, pas simplement formelle).

Holmström et al. (2023) identifient trois modèles de supervision
reconnus par l\textquotesingle AI Act : \emph{Human-in-the-loop} (HITL)
où chaque décision requiert approbation humaine,
\emph{Human-on-the-loop} (HOTL) où le système opère de manière autonome
avec surveillance et capacité d\textquotesingle intervention humaine, et
\emph{Human-in-command} (HIC) où les humains maintiennent un contrôle
stratégique sans intervenir dans les décisions opérationnelles
individuelles. Fink (2025) détaille les cinq obligations de
l\textquotesingle Article 14(4) : compréhension des capacités et
limitations, conscience du biais d\textquotesingle automation,
interprétation correcte des outputs, capacité
d\textquotesingle override, et capacité d\textquotesingle arrêt.

Pour les systèmes multi-agents stigmergiques, le mode HOTL apparaît
comme le plus adapté : les agents opèrent de manière autonome via
coordination indirecte, tandis que les superviseurs humains monitent les
artefacts partagés et interviennent quand nécessaire. Les guardrails
environnementaux (Grisold et al., 2025) incarnent le \emph{design for
oversight} en intégrant la supervision dès la conception de
l\textquotesingle environnement partagé.

\subsection{Meaningful Human Control : tracking et
tracing}\label{meaningful-human-control--tracking-et-tracing}

Santoni de Sio et van den Hoven (2018) proposent dans \emph{Frontiers in
Robotics and AI} un cadre conceptuel pour la supervision humaine
effective des systèmes autonomes complexes, structuré autour de deux
conditions nécessaires. La condition de \emph{tracking} exige que le
système réponde de manière appropriée aux raisons morales identifiées
par les humains (Santoni de Sio \& van den Hoven, 2018, p. 3),
garantissant que les décisions reflètent les considérations éthiques et
normatives humaines même lorsque le contrôle opérationnel direct est
impossible. La condition de \emph{tracing} exige qu\textquotesingle il
soit possible de remonter les résultats du système à au moins un humain
responsable (Santoni de Sio \& van den Hoven, 2018, p. 4), garantissant
l\textquotesingle accountability.

Pour les systèmes multi-agents stigmergiques, cette exigence est
particulièrement critique car la coordination indirecte via
environnement peut créer des chaînes causales opaques et les
comportements émergents résultent d\textquotesingle interactions
multiples diluant potentiellement la responsabilité. Les traces dans
l\textquotesingle environnement (commits Git, logs) doivent permettre de
reconstituer qui a configuré quelles contraintes et approuvé quelles
politiques.

\subsection{Dimensions organisationnelles : résistance et
acceptabilité}\label{dimensions-organisationnelles--ruxe9sistance-et-acceptabilituxe9}

L\textquotesingle introduction de systèmes multi-agents dans des
workflows de développement logiciel constitue une transformation
organisationnelle profonde dont la gouvernance technique ne suffit pas à
garantir le succès. Markus (1983), dans un article fondateur publié dans
\emph{Communications of the ACM}, établit que les échecs
d\textquotesingle implémentation de systèmes
d\textquotesingle information ne résultent pas principalement de
déficiences techniques, mais émergent lorsque le SI redistribue pouvoir
et contrôle au sein de l\textquotesingle organisation (Markus, 1983, p.
430). Cette théorie interactionniste, qui attribue la résistance à
l\textquotesingle interaction entre les caractéristiques du système et
le contexte organisationnel plutôt qu\textquotesingle aux seules
caractéristiques du système ou des utilisateurs, est particulièrement
pertinente pour l\textquotesingle orchestration stigmergique. En effet,
le passage d\textquotesingle une coordination hiérarchique à une
coordination émergente réalloue fondamentalement trois dimensions
organisationnelles : qui décide (les décisions de design peuvent
basculer des développeurs seniors vers des algorithmes de coordination),
qui est évalué (la performance individuelle devient visible via des
métriques automatisées de contribution aux artefacts partagés), et qui
contrôle le processus (la coordination peut passer de chefs de projets
humains à des orchestrateurs algorithmiques). Markus (1983) conclut que
la gouvernance doit inclure un travail explicite sur les enjeux de
pouvoir et la re-légitimation des rôles humains, une recommandation qui
demeure essentielle quarante ans après sa formulation.

Kim et Kankanhalli (2009) approfondissent cette analyse dans \emph{MIS
Quarterly} en mobilisant la théorie du biais de statu quo pour modéliser
la résistance des utilisateurs. Leur modèle identifie trois catégories
de coûts perçus du changement, les coûts de transition
(apprentissage, adaptation des routines), les coûts
d\textquotesingle incertitude (imprévisibilité des résultats) et les
coûts de submersion (\emph{sunk costs} dans les pratiques existantes)
, auxquels s\textquotesingle ajoutent deux biais cognitifs
amplificateurs : l\textquotesingle aversion aux pertes et
l\textquotesingle illusion de contrôle (Kim \& Kankanhalli, 2009, p.
575). Pour l\textquotesingle orchestration stigmergique, les coûts
d\textquotesingle incertitude sont particulièrement saillants car le
comportement émergent est, par définition, moins prévisible que le
comportement d\textquotesingle un workflow orchestré hiérarchiquement.
Kim et Kankanhalli (2009) recommandent que les systèmes soient gouvernés
comme des programmes de transformation organisationnelle : réduire les
coûts de basculement par des phases d\textquotesingle adoption
graduelles, rendre la valeur tangible rapidement (\emph{quick wins}), et
prévoir des mécanismes de retour en arrière réduisant le risque perçu.

Lee et See (2004) proposent dans \emph{Human Factors} un modèle de
confiance calibrée en l\textquotesingle automation, qui éclaire une
dimension complémentaire de la résistance. Leur cadre distingue deux
pathologies de confiance : la sur-confiance (\emph{complacency})
conduisant à une supervision insuffisante, particulièrement
dangereuse lorsque les agents LLM produisent des hallucinations
convaincantes dans du code, et la sous-confiance (\emph{disuse})
conduisant au rejet de systèmes performants, ce qui annule les gains de
productivité attendus. La calibration de la confiance dépend de la
transparence du processus décisionnel et de la prédictibilité des
résultats (Lee \& See, 2004), deux propriétés que la stigmergie peut
favoriser : les artefacts partagés (phéromones numériques, logs de
coordination) rendent le processus de coordination observable, et les
patterns de dépôt/lecture de phéromones créent des régularités
identifiables par les opérateurs humains.

Dietvorst, Simmons et Massey (2015) documentent dans le \emph{Journal of
Experimental Psychology: General} le phénomène
d\textquotesingle{}\emph{algorithm aversion}, la tendance
systématique à rejeter les algorithmes après avoir observé même une
seule erreur, tout en tolérant des taux d\textquotesingle erreur
supérieurs chez les décideurs humains. Ce biais est particulièrement
pertinent pour les systèmes multi-agents où les erreurs sont visibles
(un commit incorrect, une transformation de code erronée) et
attribuables au système algorithmique. Dans une étude complémentaire,
Dietvorst et al. (2018) démontrent dans \emph{Management Science} que
permettre aux utilisateurs de modifier même légèrement les outputs
algorithmiques, par exemple ajuster un paramètre de 2\%, réduit
significativement l\textquotesingle aversion, même lorsque cette
modification n\textquotesingle améliore pas objectivement les résultats.
Cette découverte a des implications directes pour la conception des
guardrails stigmergiques : offrir aux développeurs la possibilité de
paramétrer les seuils de phéromones, les critères de validation
automatique et les niveaux d\textquotesingle autonomie des agents peut
considérablement améliorer l\textquotesingle acceptabilité du système.

Turel et Kalhan (2023), dans \emph{MIS Quarterly}, approfondissent
encore cette analyse en montrant que l\textquotesingle aversion
algorithmique fonctionne comme un préjugé implicite mesurable par
l\textquotesingle Implicit Association Test (IAT), révélant que les
individus associent automatiquement les algorithmes à des attributs
négatifs (manque de fiabilité, froideur) même lorsque leurs jugements
explicites sont neutres. Cependant, et c\textquotesingle est un résultat
encourageant pour le déploiement organisationnel, cette aversion
implicite est transitoire et diminue significativement avec
l\textquotesingle expérience d\textquotesingle utilisation (Turel \&
Kalhan, 2023, p. 1385). Ce résultat suggère que les phases pilotes et
les programmes d\textquotesingle exposition progressive sont des leviers
efficaces de gestion du changement.

García-Ruiz et Rocchi (2025) apportent une dimension éthique
complémentaire essentielle en examinant, dans \emph{Business Ethics
Quarterly}, si le travail peut rester significatif (\emph{meaningful
work}) sous management algorithmique. Leur analyse, ancrée dans la
philosophie des vertus d\textquotesingle Alasdair MacIntyre, argumente
que le travail significatif requiert l\textquotesingle exercice de
l\textquotesingle agentivité morale, la capacité de prendre des
décisions délibérées reflétant des valeurs et un jugement professionnel.
Le management algorithmique, lorsqu\textquotesingle il réduit le rôle
humain à celui de simple exécutant de prescriptions automatisées, peut
éroder cette agentivité et transformer le travail en activité vide de
sens (García-Ruiz \& Rocchi, 2025). Pour l\textquotesingle orchestration
stigmergique, cette analyse implique que la conformité à
l\textquotesingle AI Act ne suffit pas : il faut également identifier
des points de contrôle critiques où l\textquotesingle intervention
humaine n\textquotesingle est pas seulement autorisée mais véritablement
significative, c\textquotesingle est-à-dire où le jugement humain
apporte une valeur que l\textquotesingle automatisation ne peut fournir,
préservant ainsi les conditions de l\textquotesingle agentivité
professionnelle des développeurs.

Ces six perspectives convergent vers des implications intégrées pour le
déploiement d\textquotesingle orchestrateurs stigmergiques.
L\textquotesingle adoption doit être progressive et gouvernée comme un
programme de transformation organisationnelle (Markus, 1983 ; Kim \&
Kankanhalli, 2009). Les guardrails doivent être paramétrables par les
utilisateurs pour réduire l\textquotesingle aversion algorithmique
(Dietvorst et al., 2018). La transparence des mécanismes de coordination
, rendue possible par la nature même des artefacts stigmergiques,
est essentielle pour maintenir une confiance calibrée (Lee \& See,
2004). Des phases pilotes doivent être prévues pour capitaliser sur la
transience de l\textquotesingle aversion implicite (Turel \& Kalhan,
2023). Enfin, les points de contrôle humains doivent préserver
l\textquotesingle agentivité professionnelle et le sens du travail
(García-Ruiz \& Rocchi, 2025), allant au-delà d\textquotesingle une
simple supervision formelle.

\subsection{Responsabilité morale distribuée et gouvernance éthique des
systèmes
swarm}\label{responsabilituxe9-morale-distribuuxe9e-et-gouvernance-uxe9thique-des-systuxe8mes-swarm}

Floridi (2016) introduit dans \emph{Philosophical Transactions of the
Royal Society A} le concept de \emph{Distributed Moral Actions} (DMAs)
: actions moralement significatives résultant
d\textquotesingle interactions locales individuellement neutres. Ce
cadre est essentiel pour comprendre comment des comportements émergents
moralement significatifs résultent d\textquotesingle interactions
stigmergiques individuellement neutres.

Winfield et al. (2025), dans la même revue, abordent spécifiquement la
gouvernance éthique des systèmes swarm en milieu réel, observant que
dans les systèmes conventionnels, une propriété émergente est un bug qui
doit être corrigé, mais que dans les systèmes swarm,
l\textquotesingle émergence est une propriété désirée plutôt
qu\textquotesingle un bug (Winfield et al., 2025). Cette distinction est
fondamentale pour notre recherche : les comportements émergents des
orchestrations stigmergiques ne doivent pas être supprimés mais
gouvernés, ce qui nécessite des cadres de gouvernance spécifiques allant
au-delà des approches traditionnelles d\textquotesingle assurance
qualité.

Mukherjee et Chang (2025) proposent le framework
d\textquotesingle{}\emph{Operational Agency} pour tracer
l\textquotesingle intention et la responsabilité dans les systèmes
multi-agents IA, fournissant un mécanisme de reconstruction causale des
décisions collectives.

\subsection{Sécurité des systèmes
multi-agents}\label{suxe9curituxe9-des-systuxe8mes-multi-agents}

Chen et al. (2024), dans un article publié à NeurIPS 2024, documentent
AgentPoison, la première attaque backdoor ciblant les bases de
connaissances RAG et mémoires persistantes des agents LLM, avec un taux
de succès d\textquotesingle au moins 80 \% pour un taux de poison
inférieur à 0,1 \% (Chen et al., 2024). L\textquotesingle environnement
partagé, équivalent stigmergique aux phéromones numériques,
devient ainsi un vecteur d\textquotesingle attaque persistante, créant
une vulnérabilité spécifique aux architectures stigmergiques. Lee et
Tiwari (2024) documentent le \emph{Prompt Infection}, injection de
prompts LLM-to-LLM se propageant silencieusement entre agents
interconnectés, analogie virale particulièrement pertinente pour les
systèmes de coordination indirecte.

Ces vulnérabilités imposent des exigences de sécurité spécifiques pour
les orchestrations stigmergiques : authentification des traces déposées
dans l\textquotesingle environnement, validation de
l\textquotesingle intégrité des artefacts partagés, et mécanismes de
détection d\textquotesingle anomalies dans les patterns de coordination
émergents.

À ce stade, la lecture conjointe des dimensions réglementaires,
organisationnelles et sécurité montre qu\textquotesingle une
orchestration stigmergique n\textquotesingle est viable que si la
gouvernance est pensée dès le design, pas ajoutée a posteriori.

\section{Agentic BPM et workflows}\label{agentic-bpm-et-workflows}

Le domaine émergent de l\textquotesingle Agentic Business Process
Management (BPM) constitue un terrain d\textquotesingle application
naturel pour l\textquotesingle orchestration stigmergique, à
l\textquotesingle intersection de l\textquotesingle automatisation des
processus et de la coordination multi-agents.

\subsection{Définition formelle de l\textquotesingle Agentic
BPM}\label{duxe9finition-formelle-de-lagentic-bpm}

Vu et al. (2026), dans un article présenté au BPM 2025 Responsible BPM
Forum (Springer LNBIP vol. 565), établissent la première définition
formelle de l\textquotesingle Agentic BPM basée sur une étude
qualitative avec 22 praticiens. Leur conceptualisation bidirectionnelle
distingue deux perspectives complémentaires : \emph{Agents for BPM}, le
déploiement d\textquotesingle agents logiciels autonomes pour atteindre
des objectifs de processus métier, et \emph{BPM for Agents},
l\textquotesingle application du BPM comme discipline pour la
gouvernance des agents (Vu et al., 2026, p. 2). Cette dualité résout la
tension entre autonomie et contrôle en reconnaissant que les agents
nécessitent eux-mêmes une gouvernance structurée que le BPM peut
fournir.

Vu et al. (2026) formulent six recommandations opérationnelles : définir
des objectifs clairs plutôt que des instructions procédurales rigides
(R1), établir des guardrails éthiques (R2), concevoir la collaboration
humain-agent plutôt que la substitution (R3), personnaliser le
comportement des agents (R4), gérer les risques proactivement (R5), et
assurer une intégration sûre avec traçabilité (R6). La recommandation R1
reflète le principe stigmergique de coordination indirecte : les agents
réagissent aux configurations environnementales plutôt que de suivre des
scripts rigides (Heylighen, 2016a). La recommandation R6 trouve une
implémentation naturelle dans les systèmes de contrôle de version
identifiés par Bolici et al. (2016) comme médiums stigmergiques.

\subsection{Agentic Business Process Management
Systems}\label{agentic-business-process-management-systems}

Dumas et al. (2026), dans un position paper présenté au AI for BPM
Workshop 2025, définissent les Agentic BPMS (A-BPMS) comme une classe de
systèmes d\textquotesingle information \emph{process-aware} utilisant la
technologie AI agentique pour exécuter des processus métier (Dumas et
al., 2026, p. 2). Ils identifient trois caractéristiques distinctives :
les flux d\textquotesingle exécution ne sont pas entièrement
prédéterminés, les adaptations s\textquotesingle opèrent sans
changements logiciels explicites, et les opportunités
d\textquotesingle amélioration sont découvertes de manière autonome.

Dumas et al. (2026) cartographient également la trajectoire évolutive
des systèmes de gestion de processus sur plusieurs décennies, depuis les
systèmes papier (1960-1980) jusqu\textquotesingle à
l\textquotesingle Agentic AI (2025+), en passant par les SOP, les
business rules engines et le RPA. Cette progression révèle une tendance
constante vers davantage d\textquotesingle autonomie décisionnelle, les
A-BPMS représentant le point d\textquotesingle inflexion où les systèmes
peuvent raisonner de manière véritablement autonome.

\subsection{Du RPA à l\textquotesingle automatisation
cognitive}\label{du-rpa-uxe0-lautomatisation-cognitive}

L\textquotesingle évolution du RPA vers l\textquotesingle automatisation
cognitive illustre la transition progressive vers les agents autonomes
et contextualise l\textquotesingle émergence de
l\textquotesingle Agentic BPM dans une trajectoire technologique plus
longue. Le Cognitive RPA Framework (2025), publié dans
l\textquotesingle{}\emph{European Journal of Computer Science and
Information Technology}, propose l\textquotesingle hybridation de
l\textquotesingle IA avec le RPA pour créer des systèmes capables de
traiter des cas non structurés, d\textquotesingle apprendre de
l\textquotesingle expérience, et de s\textquotesingle adapter aux
variations processuelles. Cette évolution s\textquotesingle inscrit dans
la trajectoire identifiée par Dumas et al. (2026) et prépare le terrain
pour les agents véritablement autonomes.

La distinction entre RPA traditionnel et automatisation cognitive est
fondamentale pour comprendre la valeur ajoutée d\textquotesingle une
approche stigmergique. Le RPA classique opère sur des règles
déterministes appliquées à des interfaces utilisateur, il reproduit
des séquences d\textquotesingle actions prédéfinies sans comprendre le
contexte ni s\textquotesingle adapter aux variations.
L\textquotesingle automatisation cognitive, en revanche, introduit la
capacité de perception contextuelle, de raisonnement et
d\textquotesingle apprentissage, trois propriétés que les agents LLM
possèdent nativement. Le passage à l\textquotesingle Agentic BPM ajoute
une quatrième propriété : l\textquotesingle autonomie décisionnelle,
c\textquotesingle est-à-dire la capacité de l\textquotesingle agent à
choisir ses actions en fonction de l\textquotesingle état de
l\textquotesingle environnement plutôt que de suivre un script
prédéterminé. C\textquotesingle est précisément cette propriété qui rend
la coordination stigmergique pertinente : des agents autonomes
nécessitent un mécanisme de coordination qui ne repose pas sur une
orchestration centralisée prescrivant chaque action, mais sur un
environnement partagé guidant le comportement collectif.

\subsection{Systèmes adaptatifs complexes appliqués au
BPM}\label{systuxe8mes-adaptatifs-complexes-appliquuxe9s-au-bpm}

Dooley (1997) propose un modèle de changement organisationnel basé sur
les systèmes adaptatifs complexes, démontrant que les organisations
opèrent comme des CAS avec des propriétés d\textquotesingle émergence,
d\textquotesingle auto-organisation et de co-évolution avec leur
environnement. Son cadre identifie quatre caractéristiques des CAS
organisationnels, la multiplicité d\textquotesingle agents
hétérogènes, les interactions locales non linéaires, la sensibilité aux
conditions initiales et l\textquotesingle émergence de patterns macro à
partir de comportements micro, qui s\textquotesingle appliquent
directement aux systèmes multi-agents d\textquotesingle orchestration
logicielle.

Uluhan et Aydin (2014) appliquent spécifiquement cette théorie au BPM,
argumentant que les processus métier dans les organisations complexes se
comportent comme des systèmes adaptatifs nécessitant des approches de
gestion tenant compte de l\textquotesingle émergence et de
l\textquotesingle auto-organisation plutôt que d\textquotesingle une
planification rigide et linéaire. Leur contribution principale réside
dans l\textquotesingle identification de trois limites des approches BPM
traditionnelles face à la complexité organisationnelle :
l\textquotesingle hypothèse de prévisibilité (les processus suivent des
chemins prédéterminés), l\textquotesingle hypothèse de décomposabilité
(les processus complexes peuvent être décomposés en sous-processus
indépendants), et l\textquotesingle hypothèse de stabilité (les
processus changent lentement et de manière contrôlée). Les trois
hypothèses sont violées dans les contextes de transformation de code à
grande échelle, où les interdépendances entre fichiers créent des effets
de cascade imprévisibles, les refactorings parallèles interfèrent
mutuellement, et l\textquotesingle environnement technique évolue
continûment (mises à jour de dépendances, changements
d\textquotesingle API).

La convergence entre CAS et BPM stigmergique est structurelle. Dans les
deux cas, la coordination émerge des interactions locales plutôt que
d\textquotesingle une planification globale, les agents
s\textquotesingle adaptent aux changements environnementaux sans
reconfiguration centralisée, et les patterns d\textquotesingle activité
collective émergent de manière organique plutôt que prescrite. La
stigmergie fournit le mécanisme opérationnel concret qui permet aux
principes abstraits des CAS de s\textquotesingle instancier dans des
systèmes informatiques réels.

\subsection{Process Mining et LLM pour le cycle de vie
BPM}\label{process-mining-et-llm-pour-le-cycle-de-vie-bpm}

Berti et al. (2024) examinent comment le Process Mining peut être
repensé à l\textquotesingle ère des agents IA, proposant que les agents
puissent interagir directement avec les \emph{event logs} pour
découvrir, analyser et optimiser les processus de manière autonome.
Cette vision est particulièrement pertinente pour
l\textquotesingle orchestration stigmergique, où les traces
environnementales (commits, modifications de fichiers, résultats de
tests) constituent naturellement un \emph{event log} exploitable par les
techniques de Process Mining. L\textquotesingle intégration du Process
Mining dans une architecture stigmergique permettrait de boucler la
boucle d\textquotesingle apprentissage : les agents génèrent des traces
via leurs actions sur l\textquotesingle environnement, le Process Mining
analyse ces traces pour identifier des patterns
d\textquotesingle efficacité ou des goulots
d\textquotesingle étranglement, et ces analyses alimentent à leur tour
les mécanismes de phéromones (renforcement des traces associées à des
patterns efficaces, évaporation de celles associées à des impasses).

Vidgof, Bachhofner et Mendling (2023) identifient les opportunités et
défis de l\textquotesingle application des LLM à
l\textquotesingle ensemble du cycle de vie BPM, découverte de
processus, analyse de conformité, identification de goulots
d\textquotesingle étranglement, et recommandation
d\textquotesingle améliorations, tout en soulignant les risques de
hallucination et d\textquotesingle incohérence qui nécessitent des
mécanismes de validation robustes, précisément le type de guardrails
environnementaux qu\textquotesingle un cadre stigmergique permet
d\textquotesingle implémenter.

Ce détour par le BPM montre qu\textquotesingle une coordination
stigmergique peut s\textquotesingle inscrire dans un cadre de processus
déjà connu des organisations, ce qui facilite
l\textquotesingle appropriation côté métier et gouvernance.

\section{Benchmarks et évaluation des systèmes
multi-agents}\label{benchmarks-et-uxe9valuation-des-systuxe8mes-multi-agents}

L\textquotesingle étape suivante consiste à évaluer cette proposition
avec des critères robustes, au-delà de la seule performance brute.

\subsection{MultiAgentBench : évaluation de la collaboration et
compétition}\label{multiagentbench--uxe9valuation-de-la-collaboration-et-compuxe9tition}

Zhu et al. (2025) introduisent le premier benchmark systématique pour
évaluer les capacités de collaboration et de compétition dans les MAS
basés sur LLM, observant que les benchmarks existants se concentrent sur
les capacités d\textquotesingle agents individuels dans des
environnements statiques, négligeant les dynamiques
d\textquotesingle interaction (Zhu et al., 2025, p. 1). MultiAgentBench
structure l\textquotesingle évaluation autour de scénarios
collaboratifs, compétitifs et mixtes, et teste systématiquement
plusieurs topologies de coordination (étoile, chaîne, arbre, graphe). La
topologie impacte significativement la performance : les architectures
en étoile excellent pour la centralisation de décision tandis que les
graphes complets favorisent l\textquotesingle émergence de consensus
distribués (Zhu et al., 2025, p. 8). Le benchmark introduit des
métriques multidimensionnelles incluant des KPIs basés sur des jalons,
des scores de complétion, et des métriques spécifiques de qualité de
coordination.

\subsection{REALM-Bench : planification
dynamique}\label{realm-bench--planification-dynamique}

Geng et Chang (2025) introduisent REALM-Bench, évaluant les MAS dans des
contextes de planification et scheduling sous contraintes dynamiques.
Identifiant une lacune critique, la plupart des benchmarks évaluent
dans des environnements statiques alors que les déploiements réels
impliquent perturbations continues (Geng \& Chang, 2025, p. 2),
REALM-Bench intègre perturbations dynamiques, contraintes de ressources
évolutives et nécessité de réoptimisation continue. Son intégration
prête à l\textquotesingle emploi avec cinq frameworks majeurs
(LangGraph, OpenAI Swarm, CrewAI, AutoGen, ALAS) permet des comparaisons
directes standardisées (Geng \& Chang, 2025, p. 9).

\subsection{Métriques pour la génération de
code}\label{muxe9triques-pour-la-guxe9nuxe9ration-de-code}

Ghosh Paul et al. (2024) produisent la synthèse la plus complète des
métriques d\textquotesingle évaluation pour la génération de code par
LLM, classifiant les métriques en quatre catégories : similarité
textuelle (BLEU, ROUGE, CodeBLEU), correction fonctionnelle (Pass@k),
qualité logicielle (lisibilité, maintenabilité, sécurité), et métriques
sémantiques évaluant la préservation de la signification. Le benchmark
RACE (2025) propose d\textquotesingle aller au-delà de la correction
pour évaluer la lisibilité, l\textquotesingle efficacité, la
maintenabilité et la documentation du code généré.

SWE-bench (Jimenez et al., 2024) s\textquotesingle est imposé comme
référence pour l\textquotesingle évaluation des agents de génie
logiciel, structuré autour d\textquotesingle issues GitHub réelles
nécessitant compréhension du contexte, localisation de bugs et
génération de patches. SWE-bench Pro (Scale AI, 2025) étend cette
évaluation avec des problèmes plus complexes nécessitant modifications
multi-fichiers.

Deux benchmarks complémentaires sont particulièrement pertinents pour
la validation du framework proposé dans ce mémoire. TravelPlanner
(Xie et al., 2024) évalue la capacité des systèmes agents à planifier
des itinéraires sous contraintes multiples (budget, durée, transport,
hébergement), servant de benchmark généraliste pour
l\textquotesingle orchestration multi-agents au-delà du génie logiciel.
Son intérêt réside dans l\textquotesingle évaluation de la satisfaction
de contraintes dures et souples (\emph{hard/soft constraints}), une
métrique transposable à la transformation de code où les contraintes
de compilation, de tests et de conventions constituent des contraintes
analogues. Amazon PolyMigration (Amazon, 2025) constitue un benchmark
industriel dédié à la migration de code cross-langage, couvrant des
scénarios de migration réalistes à l\textquotesingle échelle de
projets complets. Son orientation industrielle en fait un complément
naturel aux benchmarks académiques pour la validation de la pertinence
pratique du framework.

\subsection{Coût-efficacité et
reproductibilité}\label{couxfbt-efficacituxe9-et-reproductibilituxe9}

Kapoor et al. (2024) établissent la nécessité de frontières de Pareto
coût-précision comme méthodologie d\textquotesingle évaluation,
argumentant que les agents qui importent sont ceux qui offrent le
meilleur compromis coût-performance plutôt que la précision maximale à
tout prix. Desai et al. (2025) examinent la reproductibilité de la
recherche sur l\textquotesingle IA, proposant des métriques
distributionnelles plutôt que déterministes pour les systèmes
intrinsèquement stochastiques.

Avec ces repères d\textquotesingle évaluation, le passage au cadre
conceptuel devient plus solide : on ne propose pas seulement une idée
d\textquotesingle architecture, on prépare aussi la manière de la tester
de façon crédible.

\section{Cadre conceptuel et identification du
gap}\label{cadre-conceptuel-et-identification-du-gap}

\subsection{Synthèse croisée de la
littérature}\label{synthuxe8se-croisuxe9e-de-la-littuxe9rature}

L\textquotesingle analyse des sept blocs thématiques révèle un paysage
de recherche riche mais fragmenté. La stigmergie est bien théorisée
comme mécanisme de coordination universel (Heylighen, 2016a, 2016b) et
formalisée pour les agents cognitifs (Ricci et al., 2007). Un mouvement de convergence se dessine d'ailleurs : des frameworks récents inspirés de l'intelligence en essaim, ainsi que des sorties industrielles (dont Kimi 2.5), adoptent progressivement cette logique stigmergique.
Les systèmes multi-agents LLM font
l\textquotesingle objet d\textquotesingle une recherche active, mais
privilégient encore majoritairement les architectures hiérarchiques dont
les limitations sont documentées (Cemri et al., 2025 ; Gao et al., 2025
; Kapoor et al., 2024). Autrement dit, les briques techniques existent
désormais en grande partie, et le défi se déplace vers
l\textquotesingle intégration et l\textquotesingle orchestration
cohérente de ces éléments. La migration de code avec LLM est validée industriellement
(Ziftci et al., 2025) mais manque de cadres
d\textquotesingle orchestration multi-agents. La gouvernance de
l\textquotesingle IA est traitée au niveau réglementaire (Fink, 2025) et
organisationnel (Grisold et al., 2025 ; Jarrahi \& Ritala, 2025) mais
sans considérer les spécificités des systèmes à coordination émergente.

\subsection{Gap identifié}\label{gap-identifiuxe9}

Le gap principal réside dans l\textquotesingle absence
d\textquotesingle un cadre intégrateur combinant :

\begin{enumerate}
\def\labelenumi{\arabic{enumi}.}
\tightlist
\item
  Les principes de coordination stigmergique (Heylighen, 2016a ; Ricci
  et al., 2007) appliqués aux agents LLM plutôt qu\textquotesingle aux
  agents réactifs simples, et positionnés comme mécanisme de coordination
  organisationnelle au sens de Malone et Crowston (1994).
\item
  Les exigences de gouvernance et auditabilité (Grisold et al., 2025 ;
  Fink, 2025 ; Santoni de Sio \& van den Hoven, 2018) opérationnalisées
  pour des systèmes à coordination émergente.
\item
  Un ancrage managérial articulant capacités dynamiques (Teece, 2007),
  transformation des routines organisationnelles (Feldman \& Pentland,
  2003) et affordances technologiques (Strong et al., 2014) pour
  dépasser la seule perspective technique.
\item
  La validation empirique sur des benchmarks reconnus (TravelPlanner
  pour la validation généraliste, PolyMigration et SWE-bench pour la
  transformation de code) avec des métriques multidimensionnelles
  (performance, coût, gouvernabilité, conformité réglementaire).
\end{enumerate}

Aucune étude existante ne propose, à notre connaissance, un framework
d\textquotesingle orchestration multi-agents LLM à coordination
stigmergique (a) fondé théoriquement dans les sciences de gestion
(coordination organisationnelle, capacités dynamiques), (b) évalué
empiriquement sur des benchmarks reconnus avec des métriques intégrant
performance, coût, gouvernabilité et conformité réglementaire, et
(c) articulé sous forme de principes de conception
généralisables. C\textquotesingle est ce gap que le présent
mémoire vise à combler par une démarche de Design Science Research.

\subsection{Relations conceptuelles entre les blocs
thématiques}\label{relations-conceptuelles-entre-les-blocs-thuxe9matiques}

L\textquotesingle articulation entre les blocs thématiques de cette
revue révèle un réseau de relations conceptuelles structurant la
contribution de ce mémoire. Les fondements théoriques de la stigmergie
fournissent les principes de décentralisation, construction incrémentale
et flexibilité adaptative, ensuite opérationnalisés dans les mécanismes
computationnels de coordination par artefacts. L\textquotesingle ancrage
managérial (théorie de la coordination, capacités dynamiques, routines
organisationnelles, affordances technologiques) repositionne la
stigmergie comme mécanisme de coordination organisationnelle et requalifie
la transformation logicielle comme capacité dynamique stratégique.
Les contre-arguments
(Agentless, MAST, coût-efficacité) contraignent le design en fixant les
baselines à dépasser et les risques à éviter. Les cas de migration
(Google, Cursor, COBOL, évolution d\textquotesingle API) définissent le
domaine applicatif. Les analyses de gouvernance (guardrails,
principal-agent, AI Act, Meaningful Human Control) apportent les
contraintes de traçabilité, conformité et acceptabilité.
L\textquotesingle Agentic BPM fournit le cadre processuel intégrateur,
et les benchmarks (TravelPlanner, PolyMigration, MultiAgentBench,
REALM-Bench, RACE, SWE-bench) structurent la méthode
d\textquotesingle évaluation.

Cette structure révèle que le gap identifié n\textquotesingle est pas
simplement une absence de travail sur un sujet donné, mais une absence
d\textquotesingle intégration entre des champs de recherche riches mais
cloisonnés. La contribution originale de ce mémoire réside précisément
dans cette intégration, proposant un cadre unifié où les mécanismes
stigmergiques servent simultanément de support de coordination
technique, de vecteur de traçabilité pour la gouvernance, et
d\textquotesingle infrastructure d\textquotesingle évaluation pour la
validation empirique.

\subsection{Complexity Leadership Theory et transformation du rôle
managérial}\label{complexity-leadership-theory-et-transformation-du-ruxf4le-managuxe9rial}

L\textquotesingle adoption d\textquotesingle orchestrations
stigmergiques dans les organisations implique une transformation
profonde du rôle managérial. Uhl-Bien, Marion et McKelvey (2007), dans
un article fondamental publié dans \emph{The Leadership Quarterly},
proposent la Complexity Leadership Theory (CLT) comme cadre pour
comprendre le passage du leadership bureaucratique industriel au
leadership adaptatif de l\textquotesingle ère de la connaissance. La CLT
distingue trois fonctions de leadership interdépendantes : le leadership
administratif (structures, planification, allocation de ressources), le
leadership adaptatif (émergence créative, apprentissage organisationnel,
innovation), et le leadership enabling (création des conditions
permettant l\textquotesingle émergence adaptive) (Uhl-Bien et al.,
2007).

Cette théorie s\textquotesingle applique directement à la gouvernance
des systèmes stigmergiques : le leadership administratif définit les
guardrails et contraintes environnementales (Grisold et al., 2025), le
leadership adaptatif reconnaît et exploite les comportements émergents
plutôt que de les supprimer, et le leadership enabling crée les
conditions (infrastructure technique, culture organisationnelle,
formation) permettant la coordination stigmergique efficace. La CLT
légitime ainsi théoriquement le passage d\textquotesingle un management
directif des workflows (orchestration hiérarchique) à un management
facilitateur de l\textquotesingle auto-organisation (orchestration
stigmergique).

Shrestha, Ben-Menahem et von Krogh (2019) complètent cette perspective
dans \emph{California Management Review} en proposant un framework des
structures décisionnelles humain-IA : délégation complète à
l\textquotesingle IA, décision séquentielle hybride, et décision agrégée
humain-IA. Ces trois structures décisionnelles correspondent à des
niveaux d\textquotesingle autonomie agentique croissants et doivent être
sélectionnées en fonction de cinq facteurs contingents incluant la
prévisibilité de la tâche, la disponibilité des données et les enjeux
éthiques (Shrestha et al., 2019). Pour l\textquotesingle orchestration
stigmergique de la migration de code, la décision séquentielle hybride
semble la plus adaptée : les agents génèrent et proposent des
transformations (première phase), les guardrails environnementaux
filtrent et valident automatiquement (deuxième phase), et les humains
interviennent pour les cas complexes ou critiques (troisième phase).

\subsection{Cadre conceptuel proposé}\label{cadre-conceptuel-proposuxe9}

Le cadre conceptuel de ce mémoire repose sur trois piliers articulant
les contributions de la littérature, structurés selon la logique du
Design Science Research (Hevner et al., 2004) qui exige un ancrage
simultané dans l\textquotesingle environnement organisationnel
(pertinence pratique) et dans la base de connaissances théoriques
(rigueur scientifique).

Pilier managérial (primaire), Coordination organisationnelle et
capacités dynamiques. La théorie de la coordination (Malone \& Crowston,
1994) positionne la stigmergie comme un mécanisme de coordination
organisationnelle distinct des hiérarchies et des marchés, gérant les
dépendances par modification environnementale. Les capacités dynamiques
(Teece, 2007) requalifient la transformation logicielle continue comme
une capacité organisationnelle stratégique que le framework
opérationnalise. Les routines organisationnelles (Feldman \& Pentland,
2003) éclairent la nature du code comme incarnation de routines
organisationnelles dont la transformation exige traçabilité et
préservation du savoir tacite. Les affordances technologiques (Strong
et al., 2014) fournissent le vocabulaire analytique pour évaluer la
valeur organisationnelle du framework. Les guardrails de Grisold et al.
(2025) structurent les contraintes environnementales. La perspective
principal-agent de Jarrahi et Ritala (2025) cadre la relation
utilisateur-agents. La Complexity Leadership Theory de Uhl-Bien et al.
(2007) théorise le passage du command-and-control à
l\textquotesingle orchestration d\textquotesingle écologies
d\textquotesingle agents. Les structures décisionnelles de Shrestha et
al. (2019) définissent les modes de collaboration humain-IA.

Pilier conceptuel (fondation), Stigmergie cognitive et coordination
par artefacts. La stigmergie cognitive (Ricci et al., 2007) fournit le
cadre théorique pour des agents LLM rationnels se coordonnant via des
artefacts numériques manipulables. Le modèle de coordination de
Heylighen (2016a, 2016b) guide la conception des mécanismes de dépôt,
évaporation et renforcement des traces. Les espaces de tuples de
Gelernter (1985) et leur extension distribuée TOTA (Mamei \& Zambonelli,
2009) fournissent l\textquotesingle infrastructure computationnelle. Les
systèmes adaptatifs complexes (Holland, 1995 ; Kauffman, 1993 ; Dooley,
1997) fondent théoriquement l\textquotesingle émergence de
l\textquotesingle ordre à partir d\textquotesingle interactions locales.

Pilier évaluation et conformité, Rigueur empirique et gouvernance.
Le Meaningful Human Control (Santoni de Sio \& van den Hoven, 2018)
opérationnalise les exigences de l\textquotesingle Article 14 de
l\textquotesingle EU AI Act (Fink, 2025). Les frontières de Pareto
coût-précision de Kapoor et al. (2024) structurent
l\textquotesingle évaluation. Le baseline Agentless (Xia et al., 2024)
et le diagnostic MAST (Cemri et al., 2025) fournissent les références
comparatives. TravelPlanner (Xie et al., 2024) et Amazon PolyMigration
(Amazon, 2025) complètent le dispositif d\textquotesingle évaluation
empirique, le premier comme benchmark généraliste de planification sous
contraintes, le second comme benchmark industriel de migration
cross-langage.

\section{Design de recherche}\label{design-de-recherche}

\subsection{Approche méthodologique : Design Science Research}\label{approche-muxe9thodologique--design-science-research}

Cette recherche adopte une démarche de \emph{Design Science Research}
(DSR), méthodologie établie en systèmes d\textquotesingle information par
Hevner, March, Park et Ram (2004) dans \emph{MIS Quarterly} et
opérationnalisée par le modèle de processus DSRM de Peffers, Tuunanen,
Rothenberger et Chatterjee (2007). Le DSR est la méthodologie appropriée
pour cette recherche car l\textquotesingle objectif n\textquotesingle est
pas de tester des hypothèses causales (approche positiviste) ni de
comprendre des significations subjectives (approche interprétiviste),
mais de \emph{concevoir et évaluer un artefact} résolvant un problème
organisationnel identifié (Hevner et al., 2004, p. 77).

L\textquotesingle artefact central est un framework
d\textquotesingle orchestration stigmergique pour systèmes multi-agents
LLM, conçu d\textquotesingle abord comme architecture généraliste puis
spécialisé pour la transformation de code. Suivant la taxonomie de
Hevner et al. (2004), cet artefact comprend trois niveaux :
un \emph{modèle} (l\textquotesingle architecture stigmergique définissant
types de marqueurs, règles comportementales des agents et propriétés
de coordination émergente), une \emph{instanciation} (le prototype
fonctionnel validant la faisabilité), et des \emph{principes de
conception} généralisables extraits des
cycles de conception-évaluation (Gregor, Chandra Kruse \& Seidel, 2020).

Le positionnement de la contribution suit le cadre de Gregor et Hevner
(2013) : il s\textquotesingle agit d\textquotesingle une exaptation
au sens de Gregor et Hevner,
c\textquotesingle est-à-dire l\textquotesingle application d\textquotesingle un
concept de solution mature (la stigmergie, étudiée depuis 1959) à un
domaine problématique nouveau (l\textquotesingle orchestration
d\textquotesingle agents LLM). Ce positionnement maximise la
revendication de nouveauté tout en s\textquotesingle appuyant sur des
fondations théoriques solides.

\subsection{Processus DSRM}\label{processus-dsrm}

Le processus de recherche suit les six activités du DSRM (Peffers et al.,
2007), avec un point d\textquotesingle entrée centré sur le problème :

\emph{Activité 1, Identification et motivation du problème.}
Les organisations investissent massivement dans la modernisation logicielle,
mais les approches actuelles échouent à l\textquotesingle échelle. Les
architectures multi-agents hiérarchiques présentent des taux
d\textquotesingle échec de 41\% à 86,7\% (Cemri et al., 2025), des surcoûts
en tokens de 4x à 220x (Gao et al., 2025), et des gains marginaux
décroissants avec les modèles frontières. Simultanément, les organisations
font face à des choix stratégiques complexes, casser des monolithes,
migrer vers des microservices, moderniser incrémentalement, qui
exigent des cadres d\textquotesingle orchestration flexibles plutôt que
des workflows rigides.

\emph{Activité 2, Définition des objectifs de la solution.}
Le framework doit satisfaire cinq objectifs mesurables correspondant aux
cinq objectifs de conception : (O1) proposer une architecture
généraliste de coordination stigmergique applicable à des tâches variées ;
(O2) démontrer l\textquotesingle émergence d\textquotesingle une
spécialisation des agents sans assignation explicite de rôles ;
(O3) surpasser les frameworks centralisés sur le benchmark TravelPlanner
en score final et satisfaction des contraintes dures ;
(O4) adapter le framework à la transformation de code et atteindre des
résultats compétitifs sur PolyMigration et SWE-bench ;
(O5) produire une trace d\textquotesingle audit complète satisfaisant
les exigences de traçabilité et de conformité réglementaire, validée
par un panel d\textquotesingle experts.

\emph{Activité 3, Conception et développement.}
Le prototype sera développé en trois itérations documentées, chaque
itération étant informée par les résultats de la précédente (conformément
au principe de recherche itérative (sixième ligne directrice de Hevner
et al., 2004). Itération 1 : test des rôles agents de base et des
marqueurs stigmergiques sur des tâches simples (TravelPlanner, benchmark
de planification généraliste). Itération 2 : raffinement de la
sémantique des marqueurs et des mécanismes de résolution de conflits sur
des migrations de complexité modérée (Amazon PolyMigration). Itération
3 : introduction de comportements émergents sophistiqués
(spécialisation, priorisation par intensité de phéromone) sur des
transformations multi-fichiers (SWE-bench). Chaque décision de
conception sera tracée à une théorie fondatrice, comme
l\textquotesingle exige la rigueur DSR.

\emph{Activité 4, Démonstration.}
Application du prototype à 2-3 scénarios concrets de transformation de
code de complexité croissante sur des projets open source, permettant
la reproductibilité.

\emph{Activité 5, Évaluation.}
L\textquotesingle évaluation suit le framework FEDS (Venable, Pries-Heje
\& Baskerville, 2016) avec une stratégie centrée sur les risques
techniques et l\textquotesingle efficacité, en trois épisodes :
(1) Benchmarks techniques (TravelPlanner pour la validation généraliste,
PolyMigration et SWE-bench pour la validation spécialisée) comparant
stigmergique vs. single-agent vs. hiérarchique sur les métriques de
précision, coût, temps et \emph{overhead} de coordination ;
(2) Étude de cas sur un projet open source réaliste documentant les
comportements émergents et les patterns de coordination ;
(3) Évaluation par experts (5-8 professionnels, architectes logiciels,
managers SI, ingénieurs IA) évaluant utilité, faisabilité,
gouvernabilité et valeur organisationnelle du framework.

\emph{Activité 6, Communication.}
Le présent mémoire constitue le vecteur principal de communication.
Un article ciblant un journal classé ABS/FNEGE est envisagé en
extension.

\subsection{Conformité aux sept lignes directrices de Hevner et al. (2004)}\label{conformituxe9-aux-sept-guidelines}

Le tableau suivant synthétise la conformité du design de recherche
aux sept guidelines du DSR :

\begin{longtable}[]{@{}p{0.28\linewidth}p{0.65\linewidth}@{}}
\toprule
\textbf{Ligne directrice (Hevner et al., 2004)} & \textbf{Application dans ce mémoire} \\
\midrule
\endhead
G1. Conception d\textquotesingle un artefact & Framework stigmergique (modèle +
instanciation + principes de conception) \\
G2. Pertinence du problème & Échecs documentés des MAS hiérarchiques ;
besoins industriels de modernisation flexible \\
G3. Évaluation rigoureuse & Trois épisodes FEDS : benchmarks quantitatifs,
étude de cas, évaluation par experts \\
G4. Contributions à la recherche & Principes de conception généralisables ;
extension de la théorie de la coordination aux agents LLM \\
G5. Rigueur scientifique & Ancrage théorique multi-disciplinaire ;
métriques standardisées ; comparaison aux baselines \\
G6. Processus itératif de recherche & Trois itérations documentées avec
apprentissage inter-cycles \\
G7. Communication des résultats & Double audience : technique
(réplicabilité) et managériale (implications organisationnelles) \\
\bottomrule
\end{longtable}

\subsection{Champ d\textquotesingle étude}\label{champ-duxe9tude}

Le champ d\textquotesingle étude se situe à
l\textquotesingle intersection de trois domaines : le génie logiciel
(migration et transformation de code), les systèmes multi-agents
(orchestration et coordination), et le management des systèmes
d\textquotesingle information (gouvernance, coordination
organisationnelle et conformité). Le terrain empirique est constitué de
benchmarks reconnus (TravelPlanner, PolyMigration, SWE-bench) et de
dépôts de code open source permettant l\textquotesingle évaluation
reproductible. Le framework est conçu comme une architecture
généraliste, validée d\textquotesingle abord sur un benchmark de
planification (TravelPlanner), puis spécialisée et évaluée sur des
tâches de transformation de code.
Ce positionnement interdisciplinaire est cohérent avec la formation de
l\textquotesingle EMLV qui vise à former des managers comprenant les
enjeux technologiques et capables de gouverner
l\textquotesingle innovation.

\subsection{Collecte de données}\label{collecte-de-donnuxe9es}

La collecte de données s\textquotesingle inscrit dans le cadre DSRM et
porte sur trois sources complémentaires.

Source 1, Données du processus de conception.
Chaque itération de conception produira une documentation structurée :
décisions de conception et leurs justifications théoriques (théories
fondatrices), alternatives considérées et raisons de leur abandon,
résultats intermédiaires et apprentissages. Cette documentation
constitue la matière première pour l\textquotesingle extraction des
principes de conception.

Source 2, Données d\textquotesingle évaluation quantitative.
Le prototype générera des métriques sur chaque benchmark : taux de
succès par tâche et par projet, coût en tokens par transformation,
temps d\textquotesingle exécution, taux de rollback, \emph{overhead}
de coordination, couverture de tests avant/après, et traces
d\textquotesingle exécution (séquences d\textquotesingle actions des
agents, patterns de lecture/écriture des phéromones, cas
d\textquotesingle émergence ou d\textquotesingle impasse). Les
frontières de Pareto coût-précision (Kapoor et al., 2024) permettront de
situer l\textquotesingle approche par rapport aux alternatives.

Source 3, Données d\textquotesingle évaluation par experts.
Un panel de 5-8 experts (architectes logiciels, managers SI, ingénieurs
IA) évaluera le framework via un protocole structuré : présentation du
concept stigmergique, démonstration live, questionnaire avec items
Likert sur utilité, faisabilité, gouvernabilité et valeur
organisationnelle, puis questions ouvertes. Ce dispositif remplace les
entretiens semi-directifs exploratoires initialement envisagés : dans
le cadre DSR, les experts évaluent l\textquotesingle artefact plutôt que
de fournir des données exploratoires (Venable et al., 2016). Des
entretiens complémentaires informels avec des praticiens du terrain
(data scientists, responsables IA, managers de programmes de
transformation) pourront enrichir la compréhension du contexte
organisationnel et nourrir la discussion.

\subsection{Analyse des données}\label{analyse-des-donnuxe9es}

Les données quantitatives seront analysées selon plusieurs cadres
complémentaires. Les frontières de Pareto coût-précision (Kapoor et al.,
2024) permettront de situer l\textquotesingle approche stigmergique par
rapport à l\textquotesingle ensemble des alternatives existantes. La
comparaison au baseline Agentless (Xia et al., 2024) établira si la
complexité additionnelle de la coordination multi-agents est justifiée.
La taxonomie MAST (Cemri et al., 2025) servira de grille de lecture
des modes d\textquotesingle échec observés.

Les données qualitatives (traces de conception, évaluations experts)
seront analysées pour extraire des principes de conception selon le
format de Gregor et al. (2020) : « Pour [contexte/classe
d\textquotesingle utilisateurs], [caractéristique de
l\textquotesingle artefact] produira [résultat] parce que
[mécanisme/justification] ». L\textquotesingle objectif est de formuler
4-6 principes de conception généralisables au-delà du seul cas de la
transformation de code.

\subsection{Limites anticipées}\label{limites-anticipuxe9es}

Plusieurs limites méthodologiques sont anticipées. Le prototype sera
évalué sur un ensemble de benchmarks et de cas, mais la
généralisabilité restera conditionnée par le volume et la diversité des
scénarios effectivement testés. Le panel d\textquotesingle experts
(5-8 personnes) vise la richesse qualitative mais ne prétend pas à la
représentativité statistique. Les résultats dépendront des modèles LLM
disponibles au moment du développement, dont les capacités évoluent
rapidement. Le protocole expérimental peut engendrer des coûts
d\textquotesingle API non négligeables, ce qui peut contraindre
l\textquotesingle ampleur des campagnes de test. Le contexte
mono-chercheur du mémoire limite les possibilités de triangulation par
les investigateurs. Enfin, le niveau de contribution visé est celui de
la théorie de conception émergente au sens de Gregor et Hevner (2013), c\textquotesingle est-à-dire
des principes de conception soutenus par une instanciation et des
preuves d\textquotesingle évaluation initiales, nécessitant des
recherches futures pour atteindre le statut de théorie de conception
complète. Ces limites sont reconnues et seront discutées dans
l\textquotesingle analyse des résultats.

\newpage
\section{Bibliographie}\label{bibliographie}

Amazon. (2024). \emph{Amazon Q Developer Agent: Java migration
capabilities}. \url{https://aws.amazon.com/q/developer/}

Amazon. (2025). \emph{PolyMigration: A benchmark for cross-language code
migration}. Amazon Web Services.
\url{https://aws.amazon.com/q/developer/polymigration/}

Anthropic. (2025). \emph{When to use multi-agent systems (and when not
to)}. Claude Documentation. \url{https://docs.anthropic.com/}

Anthropic. (2025b, December 2). \emph{How AI is transforming work at
Anthropic}.
\url{https://www.anthropic.com/news/how-ai-is-transforming-work-at-anthropic}

Anthropic. (2025c). \emph{Model Context Protocol (MCP)}.
\url{https://modelcontextprotocol.io/}

Anthropic. (2026). \emph{Claude Code}. \url{https://claude.ai/code}

Barke, S., James, M. B., \& Polikarpova, N. (2023). Grounded Copilot:
How programmers interact with code-generating models. \emph{Proceedings
of the ACM on Programming Languages, 7}(OOPSLA1), 85-111.
\url{https://doi.org/10.1145/3586030}

Berti, A., Maatallah, M., Jessen, U., Sroka, M., \& Ghannouchi, S. A.
(2024). Re-thinking process mining in the AI-based agents era.
\emph{arXiv preprint arXiv:2408.07720}.
\url{https://doi.org/10.48550/arXiv.2408.07720}

Bolici, F., Howison, J., \& Crowston, K. (2016). Stigmergic coordination
in FLOSS development teams: Integrating explicit and implicit
mechanisms. \emph{Cognitive Systems Research, 38}, 14-22.
\url{https://doi.org/10.1016/j.cogsys.2015.12.003}

Bonabeau, E., Dorigo, M., \& Theraulaz, G. (1999). \emph{Swarm
intelligence: From natural to artificial systems}. Oxford University
Press.

Carriero, N., \& Gelernter, D. (1992). Coordination languages and their
significance. \emph{Communications of the ACM, 35}(2), 97-107.
\url{https://doi.org/10.1145/129630.129635}

Cemri, M., Pan, M. Z., Yang, S., et al. (2025). Why do multi-agent LLM
systems fail? \emph{arXiv preprint arXiv:2503.13657}.
\url{https://doi.org/10.48550/arXiv.2503.13657}

Chari, A., Tiwari, A., Lian, R., Reddy, S., \& Zhou, B. (2025).
Pheromone-based learning of optimal reasoning paths. \emph{arXiv
preprint arXiv:2501.19278}.
\url{https://doi.org/10.48550/arXiv.2501.19278}

Chen, Z., Zhao, Z., Zhou, A., Huang, Y., Wang, D., Poon, J., \& Li, J.
(2024). AgentPoison: Red-teaming LLM agents via poisoning memory or
knowledge bases. In \emph{Advances in Neural Information Processing
Systems 37 (NeurIPS 2024)}. \url{https://arxiv.org/abs/2407.12784}

CodeRosetta. (2024). Pushing the boundaries of unsupervised code
translation for parallel programming. \emph{Advances in Neural
Information Processing Systems 37 (NeurIPS 2024)}.

Cognitive RPA Framework. (2025). Cognitive RPA: A framework for
hybridizing artificial intelligence with robotic process automation.
\emph{European Journal of Computer Science and Information Technology,
13}(19), 64-78.

Creswell, J. W., \& Creswell, J. D. (2018). \emph{Research design:
Qualitative, quantitative, and mixed methods approaches} (5th ed.). SAGE
Publications.

Cursor. (2025). Scaling long-running autonomous coding. \emph{Cursor
Blog}. \url{https://cursor.com/blog/scaling-agents}

Desai, A., et al. (2025). What is reproducibility in artificial
intelligence and machine learning research? \emph{arXiv preprint
arXiv:2407.10239}. \url{https://doi.org/10.48550/arXiv.2407.10239}

Dietvorst, B. J., Simmons, J. P., \& Massey, C. (2015). Algorithm
aversion: People erroneously avoid algorithms after seeing them err.
\emph{Journal of Experimental Psychology: General, 144}(1), 114-126.
\url{https://doi.org/10.1037/xge0000033}

Dietvorst, B. J., Simmons, J. P., \& Massey, C. (2018). Overcoming
algorithm aversion: People will use algorithms if they can (even
slightly) modify them. \emph{Management Science, 64}(3), 1155-1170.
\url{https://doi.org/10.1287/mnsc.2016.2643}

Dig, D., \& Johnson, R. (2005). The role of refactorings in API
evolution. In \emph{IEEE International Conference on Software
Maintenance (ICSM 2005)} (pp. 2-11).
\url{https://doi.org/10.1109/ICSM.2005.26}

Diggs, C., Doyle, M., Madan, A., \& Scott, E. O. (2025). Leveraging LLMs
for legacy code modernization: Evaluation of LLM-generated
documentation. In \emph{2025 IEEE/ACM International Workshop on Large
Language Models for Code (LLM4Code@ICSE 2025)}.
\url{https://doi.org/10.1109/LLM4Code66737.2025.00027}

Dooley, K. (1997). A complex adaptive systems model of organization
change. \emph{Nonlinear Dynamics, Psychology, and Life Sciences, 1}(1),
69-97. \url{https://doi.org/10.1023/A:1022375910940}

Dumas, M., Leopold, H., Mendling, J., Rinderle-Ma, S., Smedt, J. D., \&
Zimoch, M. (2026). Agentic business process management systems.
\emph{arXiv preprint arXiv:2601.18833}.
\url{https://doi.org/10.48550/arXiv.2601.18833}

Fink, M. (2025). Human oversight under Article 14 of the EU AI Act.
\emph{SSRN}. \url{https://ssrn.com/abstract=5147196}

Feldman, M. S., \& Pentland, B. T. (2003). Reconceptualizing
organizational routines as a source of flexibility and change.
\emph{Administrative Science Quarterly, 48}(1), 94-118.
\url{https://doi.org/10.2307/3556620}

Floridi, L. (2016). Faultless responsibility: On the nature and
allocation of moral responsibility for distributed moral actions.
\emph{Philosophical Transactions of the Royal Society A, 374}(2083),
20160112. \url{https://doi.org/10.1098/rsta.2016.0112}

Gao, M., Li, Y., Liu, B., et al. (2025). Single-agent or multi-agent
systems? Why not both? \emph{arXiv preprint arXiv:2505.18286}.
\url{https://doi.org/10.48550/arXiv.2505.18286}

García-Ruiz, P., \& Rocchi, M. (2025). Can work remain meaningful under
algorithmic management? A MacIntyrean analysis. \emph{Business Ethics
Quarterly, 35}(1), 1-28. \url{https://doi.org/10.1017/beq.2024.36}

Gelernter, D. (1985). Generative communication in Linda. \emph{ACM
Transactions on Programming Languages and Systems, 7}(1), 80-112.
\url{https://doi.org/10.1145/2363.2433}

Geng, L., \& Chang, E. Y. (2025). REALM-Bench: A benchmark for
evaluating multi-agent systems on real-world, dynamic planning and
scheduling tasks. \emph{arXiv preprint arXiv:2502.18836}.
\url{https://doi.org/10.48550/arXiv.2502.18836}

Ghosh Paul, D., Zhu, H., \& Bayley, I. (2024). Benchmarks and metrics
for evaluations of code generation: A critical review. \emph{arXiv
preprint arXiv:2406.12655}.
\url{https://doi.org/10.48550/arXiv.2406.12655}

Grassé, P.-P. (1959). La reconstruction du nid et les coordinations
interindividuelles chez \emph{Bellicositermes natalensis} et
\emph{Cubitermes} sp. la théorie de la stigmergie. \emph{Insectes
Sociaux, 6}(1), 41-80. \url{https://doi.org/10.1007/BF02223791}

Gregor, S., \& Hevner, A. R. (2013). Positioning and presenting design
science research for maximum impact. \emph{MIS Quarterly, 37}(2),
337-355. \url{https://doi.org/10.25300/MISQ/2013/37.2.01}

Gregor, S., Chandra Kruse, L., \& Seidel, S. (2020). The anatomy of a
design principle. \emph{Journal of the Association for Information
Systems, 21}(6), 1622-1652.
\url{https://doi.org/10.17705/1jais.00649}

Grisold, T., Berente, N., \& Seidel, S. (2025). Guardrails for human-AI
ecologies: Norm-based coordination and design for predictability.
\emph{MIS Quarterly, 49}(4), 1239-1266.
\url{https://doi.org/10.25300/MISQ/2025/18058}

Heylighen, F. (2016a). Stigmergy as a universal coordination mechanism
I: Definition and components. \emph{Cognitive Systems Research, 38},
4-13. \url{https://doi.org/10.1016/j.cogsys.2015.12.002}

Heylighen, F. (2016b). Stigmergy as a universal coordination mechanism
II: Varieties and evolution. \emph{Cognitive Systems Research, 38},
50-59. \url{https://doi.org/10.1016/j.cogsys.2015.12.007}

Holland, J. H. (1995). \emph{Hidden order: How adaptation builds
complexity}. Addison-Wesley.

Holmström, J., et al. (2023). \textquotesingle Human
oversight\textquotesingle{} in the EU artificial intelligence act: What,
when and by whom? \emph{Law, Innovation and Technology} (Taylor \&
Francis). \url{https://doi.org/10.1080/17579961.2023.2293129}

Hevner, A. R., March, S. T., Park, J., \& Ram, S. (2004). Design
science in information systems research. \emph{MIS Quarterly, 28}(1),
75-105. \url{https://doi.org/10.2307/25148625}

Hong, S., Zhuge, M., Chen, J., et al. (2024). MetaGPT: Meta programming
for a multi-agent collaborative framework. \emph{arXiv preprint
arXiv:2308.00352}. \url{https://doi.org/10.48550/arXiv.2308.00352}

IBM. (2023). \emph{IBM unveils watsonx generative AI capabilities to
accelerate mainframe application modernization}.
\url{https://newsroom.ibm.com/2023-08-22-IBM-Unveils-watsonx-Generative-AI-Capabilities-to-Accelerate-Mainframe-Application-Modernization}

Jarrahi, M. H., \& Ritala, P. (2025, July 23). Rethinking AI agents: A
principal-agent perspective. \emph{California Management Review}
(Insight).
\url{https://cmr.berkeley.edu/2025/07/rethinking-ai-agents-a-principal-agent-perspective/}

Jimenez, C. E., Yang, J., Wettig, A., et al. (2024). SWE-bench: Can
language models resolve real-world GitHub issues? \emph{arXiv preprint
arXiv:2310.06770}. \url{https://doi.org/10.48550/arXiv.2310.06770}

Kapoor, S., Stroebl, B., Siegel, Z. S., Nadgir, N., \& Narayanan, A.
(2024). AI agents that matter. \emph{arXiv preprint arXiv:2407.01502}.
\url{https://doi.org/10.48550/arXiv.2407.01502}

Kauffman, S. A. (1993). \emph{The origins of order: Self-organization
and selection in evolution}. Oxford University Press.

Kim, H.-W., \& Kankanhalli, A. (2009). Investigating user resistance to
information systems implementation: A status quo bias perspective.
\emph{MIS Quarterly, 33}(3), 567-582.
\url{https://doi.org/10.2307/20650309}

Lamothe, M., Guéhéneuc, Y.-G., \& Shang, W. (2021). A systematic review
of API evolution literature. \emph{ACM Computing Surveys, 54}(8),
171:1-171:36. \url{https://doi.org/10.1145/3450268}

Lee, D., \& Tiwari, M. (2024). Prompt injection: LLM-to-LLM prompt
injection within multi-agent systems. \emph{arXiv preprint
arXiv:2410.07283}. \url{https://arxiv.org/abs/2410.07283}

Lee, J. D., \& See, K. A. (2004). Trust in automation: Designing for
appropriate reliance. \emph{Human Factors, 46}(1), 50-80.
\url{https://doi.org/10.1518/hfes.46.1.50_30392}

Liu, X., Yu, H., Zhang, H., et al. (2024). AgentBench: Evaluating LLMs
as agents. \emph{arXiv preprint arXiv:2308.03688}.
\url{https://doi.org/10.48550/arXiv.2308.03688}

Mamei, M., \& Zambonelli, F. (2009). Programming pervasive and mobile
computing applications: The TOTA approach. \emph{ACM Transactions on
Software Engineering and Methodology, 18}(4), 1-56.
\url{https://doi.org/10.1145/1538942.1538945}

Malone, T. W., \& Crowston, K. (1994). The interdisciplinary study of
coordination. \emph{ACM Computing Surveys, 26}(1), 87-119.
\url{https://doi.org/10.1145/174666.174668}

Markus, M. L. (1983). Power, politics, and MIS implementation.
\emph{Communications of the ACM, 26}(6), 430-444.
\url{https://doi.org/10.1145/358141.358148}

Marsh, L., \& Onof, C. (2008). Stigmergic epistemology, stigmergic
cognition. \emph{Cognitive Systems Research, 9}(1-2), 136-149.
\url{https://doi.org/10.1016/j.cogsys.2007.06.009}

Miles, M. B., Huberman, A. M., \& Saldaña, J. (2020). \emph{Qualitative
data analysis: A methods sourcebook} (4th ed.). SAGE Publications.

Mukherjee, A., \& Chang, H. (2025). Operational agency: A framework for
tracing intent and liability in multi-agent artificial intelligence
systems. \emph{SSRN Electronic Journal}.
\url{https://doi.org/10.2139/ssrn.5136065}

Parunak, H. V. D. (1997). "Go to the ant": Engineering principles from
natural agent systems. \emph{Annals of Operations Research, 75}, 69-101.
\url{https://doi.org/10.1023/A:1018980001403}

OpenAI. (2026, February 5). \emph{Introducing GPT‑5.3‑Codex}.
\url{https://openai.com/index/introducing-gpt-5-3-codex/}

OpenAI. (2022, November 30). \emph{Introducing ChatGPT}.
\url{https://openai.com/index/chatgpt/}

OpenAI. (2026b). \emph{Codex CLI Documentation}.
\url{https://developers.openai.com/codex/cli}

Peng, S., Kalliamvakou, E., Cihon, P., \& Demirer, M. (2023). The impact
of AI on developer productivity: Evidence from GitHub Copilot.
\emph{arXiv preprint arXiv:2302.06590}.
\url{https://doi.org/10.48550/arXiv.2302.06590}

Peffers, K., Tuunanen, T., Rothenberger, M. A., \& Chatterjee, S.
(2007). A design science research methodology for information systems
research. \emph{Journal of Management Information Systems, 24}(3),
45-77. \url{https://doi.org/10.2753/MIS0742-1222240302}

Rahman, M. A. U., Schranz, M., \& Hayat, S. (2025). LLM-powered swarms:
A new frontier or a conceptual stretch? \emph{arXiv preprint
arXiv:2506.14496}. \url{https://doi.org/10.48550/arXiv.2506.14496}

Ricci, A., Omicini, A., Viroli, M., Gardelli, L., \& Oliva, E. (2007).
Cognitive stigmergy: Towards a framework based on agents and artifacts.
In D. Weyns, H. V. D. Parunak, \& F. Michel (Eds.), \emph{Environments
for Multi-Agent Systems III} (LNCS Vol. 4389, pp. 124-140). Springer.
\url{https://doi.org/10.1007/978-3-540-71103-2_7}

Rodriguez, R. R., Jr. (2026). Emergent coordination in multi-agent
systems via pressure fields and temporal decay. \emph{arXiv preprint
arXiv:2601.08129}. \url{https://doi.org/10.48550/arXiv.2601.08129}

Rozière, B., Lachaux, M.-A., Chanussot, L., \& Lample, G. (2020).
Unsupervised translation of programming languages. \emph{Advances in
Neural Information Processing Systems 33 (NeurIPS 2020)}, 20601-20611.
\url{https://doi.org/10.48550/arXiv.2009.08732}

Santoni de Sio, F., \& van den Hoven, J. (2018). Meaningful human
control over autonomous systems: A philosophical account.
\emph{Frontiers in Robotics and AI, 5}, Article 15.
\url{https://doi.org/10.3389/frobt.2018.00015}

Serugendo, G. D. M., Gleizes, M. P., \& Karageorgos, A. (2005).
Self-organization in multi-agent systems. \emph{The Knowledge
Engineering Review, 20}(2), 165-189.
\url{https://doi.org/10.1017/S0269888905000494}

Shrestha, Y. R., Ben-Menahem, S. M., \& von Krogh, G. (2019).
Organizational decision-making structures in the age of artificial
intelligence. \emph{California Management Review, 61}(4), 66-83.
\url{https://doi.org/10.1177/0008125619862257}

Sneed, H. M. (2010). Migrating from COBOL to Java. In \emph{IEEE
International Conference on Software Maintenance (ICSM 2010)} (pp.
237-246). \url{https://doi.org/10.1109/ICSM.2010.5609741}

Strong, D. M., Volkoff, O., Johnson, S. A., Pelletier, L. R., Tulu, B.,
Bar-On, I., Trudel, J., \& Garber, L. (2014). A theory of
organization-EHR affordance actualization. \emph{Journal of the
Association for Information Systems, 15}(2), 53-85.
\url{https://doi.org/10.17705/1jais.00353}

Teece, D. J. (2007). Explicating dynamic capabilities: The nature and
microfoundations of (sustainable) enterprise performance.
\emph{Strategic Management Journal, 28}(13), 1319-1350.
\url{https://doi.org/10.1002/smj.640}

Theraulaz, G., \& Bonabeau, E. (1999). A brief history of stigmergy.
\emph{Artificial Life, 5}(2), 97-116.
\url{https://doi.org/10.1162/106454699568700}

Turel, O., \& Kalhan, S. (2023). Prejudiced against the machine?
Implicit associations and the transience of algorithm aversion.
\emph{MIS Quarterly, 47}(4), 1369-1394.
\url{https://doi.org/10.25300/MISQ/2023/16860}

Uhl-Bien, M., Marion, R., \& McKelvey, B. (2007). Complexity leadership
theory: Shifting leadership from the industrial age to the knowledge
era. \emph{The Leadership Quarterly, 18}(4), 298-318.
\url{https://doi.org/10.1016/j.leaqua.2007.04.002}

Uluhan, E., \& Aydin, M. N. (2014). Complex adaptive systems theory in
the context of business process management. In \emph{S-BPM ONE 2014:
Scientific Research} (pp. 11-23). Springer.
\url{https://doi.org/10.1007/978-3-319-06191-7_2}

Vidgof, M., Bachhofner, S., \& Mendling, J. (2023). Large language
models for business process management: Opportunities and challenges.
\emph{arXiv preprint arXiv:2304.04309}.
\url{https://doi.org/10.48550/arXiv.2304.04309}

Vu, H., Klievtsova, N., Leopold, H., Rinderle-Ma, S., \& Kampik, T.
(2026). Agentic business process management: Practitioner perspectives
on agent governance in business processes. In \emph{BPM 2025 Responsible
BPM Forum. Lecture Notes in Business Information Processing} (vol. 565).
Springer. \url{https://doi.org/10.48550/arXiv.2504.03693}

Venable, J., Pries-Heje, J., \& Baskerville, R. (2016). FEDS: A
framework for evaluation in design science research. \emph{European
Journal of Information Systems, 25}(1), 77-89.
\url{https://doi.org/10.1057/ejis.2014.36}

Wang, X., Li, B., Song, Y., et al. (2025). OpenHands: An open platform
for AI software developers as generalist agents. \emph{arXiv preprint
arXiv:2407.16741}. \url{https://doi.org/10.48550/arXiv.2407.16741}

Weyns, D., Omicini, A., \& Odell, J. (2007). Environment as a first
class abstraction in multiagent systems. \emph{Autonomous Agents and
Multi-Agent Systems, 14}(1), 5-30.
\url{https://doi.org/10.1007/s10458-006-0012-0}

Winfield, A. F. T., et al. (2025). On the ethical governance of swarm
robotic systems in the real world. \emph{Philosophical Transactions of
the Royal Society A, 383}(2289), 20240142.
\url{https://doi.org/10.1098/rsta.2024.0142}

Xia, C. S., Deng, Y., Dunn, S., \& Zhang, L. (2024). Agentless:
Demystifying LLM-based software engineering agents. \emph{arXiv preprint
arXiv:2407.01489}. \url{https://doi.org/10.48550/arXiv.2407.01489}

Xie, J., Zhang, K., Chen, J., Zhu, T., Lou, R., Tian, Y., Xiao, Y.,
\& Su, Y. (2024). TravelPlanner: A benchmark for real-world planning
with language agents. In \emph{Proceedings of the 41st International
Conference on Machine Learning (ICML 2024)}.
\url{https://doi.org/10.48550/arXiv.2402.01622}

Yan, W. (2025, June 12). Don\textquotesingle t build multi-agents.
\emph{Cognition AI Blog}. \url{https://www.cognition-labs.com/blog}

Yang, J., Jimenez, C. E., Yao, S., et al. (2024). SWE-agent:
Agent-computer interfaces enable automated software engineering.
\emph{arXiv preprint arXiv:2405.15793}.
\url{https://doi.org/10.48550/arXiv.2405.15793}

Zhang, Y., Lin, C., Tang, S., Chen, H., Zhou, S., Ma, Y., \& Tresp, V.
(2025). SwarmAgentic: Towards fully automated agentic system generation
via swarm intelligence. \emph{arXiv preprint arXiv:2506.15672}.
\url{https://doi.org/10.48550/arXiv.2506.15672}

Zhu, K., Du, H., Hong, Z., et al. (2025). MultiAgentBench: Evaluating
the collaboration and competition of LLM agents. \emph{arXiv preprint
arXiv:2503.01935}. \url{https://doi.org/10.48550/arXiv.2503.01935}

Ziftci, C., Zheng, J., Choudhary, S., Rossbach, C., Chen, W., Li, W.,
Wang, J., \& Zhang, L. (2025). Migrating code at scale with LLMs at
Google. In \emph{Proceedings of the ACM International Conference on the
Foundations of Software Engineering (FSE 2025)}.

\newpage
\appendix
\input{annexe_ia.tex}

\end{document}